\documentclass[11pt]{report}

%to allow more levels of sections and paragraphs:
\setcounter{secnumdepth}{4}

%Packages
\usepackage[utf8]{inputenc}
\usepackage[breaklinks=true,hidelinks]{hyperref}
\usepackage{graphicx}
\usepackage[spanish, english]{babel}
\usepackage{url}
\usepackage{multirow}
\usepackage{verbatim} %comentarios
\usepackage{amsmath}
\usepackage{amssymb}
\usepackage{textcomp,gensymb}
\usepackage{colortbl}
\usepackage{graphicx} % Para escalar la tabla
\usepackage[table,xcdraw]{xcolor}
\usepackage{multirow}
\usepackage{diagbox}
\usepackage{dashrule}
\usepackage{pdflscape}
\usepackage{fancyhdr} % Para personalizar encabezados y pies
\usepackage[T1]{fontenc}

\usepackage{multicol} 

\usepackage{subfig}

\usepackage{rotating}

\usepackage{pifont}

\usepackage{relsize}
\usepackage{xcolor,colortbl}
\usepackage{wrapfig}
\usepackage{bbm}
\usepackage[labelfont=bf]{caption}
\usepackage{tabularx}
\usepackage{booktabs}

\DeclareMathOperator{\arctantwo}{arctan2}

\graphicspath{ {img/} }

\usepackage[document]{ragged2e}
\usepackage{float}
\usepackage{listings}
\usepackage{titlesec}
\usepackage[left=3cm, right=3cm, top = 2.5cm, bottom = 2.5cm, includehead, includefoot]{geometry}

\usepackage{appendix}
\usepackage{listings}
\definecolor{codeback}{HTML}{F7F7F7}
\definecolor{codekw}{HTML}{0033B3}
\definecolor{codestr}{HTML}{067D17}
\definecolor{codecom}{HTML}{8C8C8C}
\lstset{
  backgroundcolor=\color{codeback},
  basicstyle=\ttfamily\footnotesize,
  keywordstyle=\color{codekw}\bfseries,
  stringstyle=\color{codestr},
  commentstyle=\color{codecom}\itshape,
  numbers=left,
  numberstyle=\tiny\color{codecom},
  numbersep=8pt,
  frame=single,
  rulecolor=\color{codecom!40},
  breaklines=true,
  breakatwhitespace=true,
  showstringspaces=false,
  tabsize=2,
  captionpos=b,
  xleftmargin=12pt,
  framexleftmargin=12pt,
  literate=%
    {á}{{\'a}}1 {é}{{\'e}}1 {í}{{\'i}}1 {ó}{{\'o}}1 {ú}{{\'u}}1
    {Á}{{\'A}}1 {É}{{\'E}}1 {Í}{{\'I}}1 {Ó}{{\'O}}1 {Ú}{{\'U}}1
    {ñ}{{\~n}}1 {Ñ}{{\~N}}1 {ü}{{\"u}}1 {¿}{{?`}}1 {¡}{{!`}}1
}
\renewcommand{\lstlistingname}{Listado}
\usepackage{afterpage}
\usepackage{xparse}

\usepackage{algorithm}
\usepackage[noend]{algpseudocode}

\usepackage{eurosym}
\usepackage{tikz}
\usetikzlibrary{arrows.meta,positioning,fit,backgrounds,babel}

\setlength{\headheight}{32.09pt}

%Config
\setlength{\parskip}{1em}
\setlength{\parindent}{0em}

\titleformat*{\section}{\color{black}\bfseries\Large}
\titleformat{\chapter}[display]
{\normalfont\bfseries}{}{0pt}{\Huge}

% Títulos de capítulo pegados a la cabecera, sin el hueco por defecto (50 pt
% arriba y una línea vacía para el número). Igual con número y sin él (Índice,
% listados, Resumen, Glosario, Bibliografía).
\makeatletter
    \renewcommand*\@makechapterhead[1]{
        {\Huge\bfseries #1\par \nobreak\vskip 24\p@}
    }
    \renewcommand*\@makeschapterhead[1]{
        {\Huge\bfseries #1\par \nobreak\vskip 24\p@}
    }
\makeatother
 
\begin{document}

\selectlanguage{spanish}

\renewcommand{\contentsname}{Índice}
\renewcommand{\partname}{Parte} 
\renewcommand{\chaptername}{Capítulo} 

\renewcommand{\appendixname}{Anexo} 
\renewcommand{\figurename}{Figura}
\renewcommand{\tablename}{Tabla} 
\renewcommand{\listfigurename}{Índice de figuras} 
\renewcommand{\glossaryname}{Glosario}


\begin{titlepage}

\newgeometry{top=2.5cm, left=2cm, right=2cm}

\definecolor{customOrange}{HTML}{ff6500}

\pagecolor{customOrange}\afterpage{\nopagecolor}


\begin{center}

 { \huge \textbf{UNIVERSIDAD POLITÉCNICA DE MADRID} \par} 

 {\vspace{0.5cm}}
 
 {\scshape  \Large \textbf{ESCUELA TÉCNICA SUPERIOR \\ DE INGENIEROS DE TELECOMUNICACIÓN} \par} 



    {\vspace{1.5cm}}
    
    
\begin{figure}[H]
    \centering
    \includegraphics[scale=0.3]{img/logoBN3.png}
    \label{logoNegro}
\end{figure}{}
 
 
   {\vspace{1cm}}


    { \vspace{1.5cm} \Large \textbf{GRADO EN INGENIERÍA Y SISTEMAS DE DATOS}}

    { \vspace{0.5cm} \Large \textbf{TRABAJO FIN DE GRADO}}

    {\vspace{1cm}}
 
   {\large \textbf{PLATAFORMA DE CONCIENCIA CIBERSITUACIONAL: DETECCIÓN DE ATAQUES MULTI-ETAPA, GRAFO DE ATAQUE RAZONADO Y PRIORIZACIÓN EXPLICABLE DE INCIDENTES} \par} 
   
   
   
   
   
   \vspace{1.5cm}
   {\itshape \Large \textbf{Javier Saguar Antona} \par}
   
   
   
  \vspace{1cm}
  \Large \textbf{2027} \par


  
   
\end{center}


\restoregeometry
\end{titlepage}

\pagenumbering{gobble} %Quitar numeracion

\clearpage

\pagenumbering{gobble}

\LARGE \textbf{Grado en Ingeniería y Sistemas de Datos}

\vspace{0.5cm}

\LARGE \textbf{TRABAJO DE FIN DE GRADO}

\vspace{0.5cm}

\large \textbf{Título}: Plataforma de conciencia cibersituacional: detección de ataques multi-etapa, grafo de ataque razonado y priorización explicable de incidentes

\large \textbf{Autor}: Javier Saguar Antona

\large \textbf{Tutora}: Carmen Sánchez-Zas

\large \textbf{Departamento}: Ingeniería de Sistemas Telemáticos

\vspace{0.5cm}

\LARGE \textbf{MIEMBROS DEL TRIBUNAL}

\vspace{0.5cm}
\large \textbf{Presidente}: \hdashrule[0.5ex]{5cm}{1pt}{2pt}

\large \textbf{Vocal}: \hdashrule[0.5ex]{6cm}{1pt}{2pt}

\large \textbf{Secretario}: \hdashrule[0.5ex]{5.1cm}{1pt}{2pt}

\large \textbf{Suplente}: \hdashrule[0.5ex]{5.4cm}{1pt}{2pt}

\vspace{2cm}

\normalsize \textbf{Los miembros del tribunal arriba nombrados acuerdan otorgar la calificación de:}

\vspace{1cm}

\hdashrule[0.5ex]{6cm}{1pt}{2pt}

\vspace{1cm}

\large \textbf{Fecha de lectura}: \hdashrule[0.5ex]{5cm}{1pt}{2pt}

\normalsize

\clearpage




\clearpage

\begin{titlepage}

\newgeometry{top=2.5cm, left=2cm, right=2cm}


\begin{center}

 { \huge \textbf{UNIVERSIDAD POLITÉCNICA DE MADRID} \par} 

 {\vspace{0.5cm}}
 
 {\scshape  \Large \textbf{ESCUELA TÉCNICA SUPERIOR \\ DE INGENIEROS DE TELECOMUNICACIÓN} \par} 



    {\vspace{1.5cm}}
    
    
\begin{figure}[H]
    \centering
    \includegraphics[scale=0.14]{img/logoescuela.jpg}
    \label{logoColor}
\end{figure}{}
 
 
   {\vspace{1cm}}


    { \vspace{1.5cm} \Large \textbf{GRADO EN INGENIERÍA Y SISTEMAS DE DATOS}}

    { \vspace{0.5cm} \Large \textbf{TRABAJO FIN DE GRADO}}

    {\vspace{1cm}}
 
   {\large \textbf{PLATAFORMA DE CONCIENCIA CIBERSITUACIONAL: DETECCIÓN DE ATAQUES MULTI-ETAPA, GRAFO DE ATAQUE RAZONADO Y PRIORIZACIÓN EXPLICABLE DE INCIDENTES} \par} 
   
   
   
   
   
   \vspace{1.5cm}
   {\itshape \Large \textbf{Javier Saguar Antona} \par}
   
   
   
  \vspace{1cm}
  \Large \textbf{2027} \par


  
   
\end{center}


\restoregeometry
\end{titlepage}

\clearpage


% Page header and footer
\pagestyle{fancy}

\fancyhf{}
\lhead{
\begin{picture}(0,0)
    \includegraphics[scale=0.04]{img/logoescuela.jpg}
\end{picture}{}
}
\rhead{\thepage}
% End

\justifying


\chapter*{Resumen}
\addcontentsline{toc}{chapter}{Resumen}

Los centros de operaciones de seguridad (SOC) se enfrentan hoy a dos problemas que se
refuerzan mutuamente. Por un lado, el volumen de alertas generadas por las herramientas
de detección supera con holgura la capacidad de análisis humano, lo que provoca el
fenómeno conocido como \textit{fatiga de alertas}: los analistas descartan o posponen
avisos que podrían ser relevantes simplemente porque no hay tiempo material para
revisarlos todos. Por otro, esas alertas llegan \textit{sin contexto}: un aviso aislado no
indica si el comportamiento observado corresponde a una técnica documentada, si ha sido
empleada por un actor conocido, ni qué activos de la organización quedan expuestos como
consecuencia.

A este panorama se suma que los ataques dirigidos no se manifiestan como eventos puntuales,
sino como \textbf{secuencias de etapas} que se desarrollan a lo largo del tiempo:
reconocimiento, acceso inicial, ejecución, persistencia, movimiento lateral y exfiltración.
Los sistemas de detección que evalúan cada evento de forma independiente son, por
construcción, incapaces de reconocer esa estructura temporal, y resultan por tanto fáciles
de evadir mediante técnicas individualmente poco llamativas.

Este Trabajo Fin de Grado propone el diseño, la implementación y la validación de una
plataforma de datos que aborda ambos problemas de forma conjunta. El sistema ingiere
telemetría heterogénea de red y de sistema operativo, la normaliza a un esquema canónico
(OCSF), detecta secuencias de ataque multi-etapa mediante procesamiento de eventos
complejos (CEP), las contextualiza contra un grafo de conocimiento de inteligencia de
amenazas alineado con MITRE ATT\&CK, razona sobre las rutas de ataque factibles hacia los
activos críticos empleando un motor lógico basado en Datalog, y emite finalmente una
priorización de incidentes \textbf{explicable}, en la que cada puntuación va acompañada de
la descomposición de los factores que la justifican.

La arquitectura se articula sobre tres pilares: la detección multi-etapa, el grafo de
ataque razonado y la priorización explicable. El objetivo último no es sustituir a las
herramientas de detección existentes, sino transformar su salida en información accionable,
aportando \textbf{conciencia cibersituacional} y soporte a la decisión al analista.

La validación se realiza sobre conjuntos de datos públicos etiquetados con MITRE ATT\&CK,
lo que permite contrastar cuantitativamente la precisión de la detección, la reducción del
volumen de alertas y la calidad de la priorización frente a líneas base establecidas.

\vspace{1em}
\noindent\textbf{Palabras clave:} conciencia cibersituacional, detección de intrusiones,
MITRE ATT\&CK, procesamiento de eventos complejos, grafos de ataque, inteligencia de
amenazas, ingeniería de datos en tiempo real, explicabilidad.

\clearpage

\chapter*{Abstract}
\addcontentsline{toc}{chapter}{Abstract}

Security Operations Centres (SOCs) currently face two mutually reinforcing problems. First,
the volume of alerts produced by detection tooling far exceeds human analysis capacity,
leading to the well-documented phenomenon of \textit{alert fatigue}: analysts dismiss or
postpone potentially relevant warnings simply because there is no time to review them all.
Second, those alerts arrive \textit{without context}: an isolated warning does not indicate
whether the observed behaviour corresponds to a documented technique, whether it has been
used by a known threat actor, or which organisational assets are consequently at risk.

Compounding this, targeted attacks do not manifest as isolated events but as
\textbf{multi-stage sequences} unfolding over time: reconnaissance, initial access,
execution, persistence, lateral movement and exfiltration. Detection systems that evaluate
each event independently are, by construction, unable to recognise that temporal structure,
and are therefore easily evaded through individually unremarkable techniques.

This Bachelor's Thesis proposes the design, implementation and validation of a data platform
addressing both problems jointly. The system ingests heterogeneous network and host
telemetry, normalises it to a canonical schema (OCSF), detects multi-stage attack sequences
via complex event processing (CEP), contextualises them against a cyber threat intelligence
knowledge graph aligned with MITRE ATT\&CK, reasons over feasible attack paths towards
critical assets using a Datalog-based logic engine, and finally produces an
\textbf{explainable} incident prioritisation in which every score is accompanied by a
breakdown of the contributing factors.

The architecture rests on three pillars: multi-stage detection, the reasoned attack graph,
and explainable prioritisation. The ultimate goal is not to replace existing detection
tooling, but to transform its output into actionable information, providing
\textbf{cyber situational awareness} and decision support to the analyst.

Validation is carried out on public datasets labelled with MITRE ATT\&CK techniques,
enabling a quantitative assessment of detection accuracy, alert volume reduction and
prioritisation quality against established baselines.

\vspace{1em}
\noindent\textbf{Keywords:} cyber situational awareness, intrusion detection, MITRE ATT\&CK,
complex event processing, attack graphs, threat intelligence, real-time data engineering,
explainability.

\clearpage

%\setlength{\parindent}{0em}
\pagenumbering{gobble}
\thispagestyle{fancy}

\noindent% prevent the box to be shifted a bit to the right
\makebox[\textwidth][r]{% create a box, with it's content aligned to the right
  \begin{minipage}{.5\textwidth}% create a content half width
    \section*{Agradecimientos}
    \justifying
    \textit{Si quieres poner dedicatoria, descomentar en main.tex.}

  \end{minipage}
}
%\clearpage\null\newpage

\clearpage
\clearpage

\pagenumbering{gobble}

\fancypagestyle{plain}{
  \thispagestyle{fancy}
  \renewcommand{\contentsname}{Índice}
}

\tableofcontents

\clearpage %%% Quita esto si no quieres las figuras
\renewcommand{\listfigurename}{Listado de figuras}
\listoffigures

\clearpage %%% Quita esto si no quieres las tablas
\renewcommand{\listtablename}{Listado de tablas}
\listoftables

\clearpage



\clearpage
\chapter*{Glosario de siglas y términos}
\addcontentsline{toc}{chapter}{Glosario}

\begin{description}

\item[APT] \textit{Advanced Persistent Threat}. Actor con capacidad, recursos y motivación
para mantener una intrusión prolongada y sigilosa en una organización objetivo.

\item[ATT\&CK] Base de conocimiento pública de MITRE que cataloga tácticas, técnicas y
procedimientos (TTP) empleados por adversarios reales, organizados en una matriz.

\item[CEP] \textit{Complex Event Processing}. Paradigma de procesamiento que detecta
patrones formados por secuencias de eventos relacionados en el tiempo, en lugar de evaluar
eventos de forma aislada.

\item[Campaña] Conjunto de alertas y eventos que el sistema reconoce como etapas sucesivas de
un mismo ataque. Es la salida del Pilar~A.

\item[CTI] \textit{Cyber Threat Intelligence}. Inteligencia de amenazas: información
contextualizada sobre actores, campañas, herramientas e indicadores de compromiso.

\item[CVE] \textit{Common Vulnerabilities and Exposures}. Identificador público y único
asignado a una vulnerabilidad conocida.

\item[CVSS] \textit{Common Vulnerability Scoring System}. Sistema estándar de puntuación de
la severidad de una vulnerabilidad.

\item[Datalog] Lenguaje declarativo de consulta y razonamiento, subconjunto de Prolog,
empleado para derivar hechos nuevos a partir de hechos base y reglas de inferencia.

\item[DLQ] \textit{Dead Letter Queue}. Cola de rechazos donde se derivan los mensajes que
no han podido procesarse, para su inspección posterior sin pérdida de información.

\item[EPSS] \textit{Exploit Prediction Scoring System}. Estimación de la probabilidad de
que una vulnerabilidad concreta sea explotada en un plazo determinado.

\item[Event time] Marca temporal correspondiente al instante en que ocurrió realmente un
evento, por oposición al \textit{processing time} (instante en que el sistema lo procesa).

\item[Hallazgo] Clase de OCSF (\textit{Detection Finding}) que representa una conclusión del
sistema de detección. Es el formato en que se emiten las campañas.

\item[IOC] \textit{Indicator of Compromise}. Artefacto observable que evidencia una
intrusión: una dirección IP, un hash, un dominio o una ruta de fichero.

\item[Kill chain] Modelo que describe un ataque como una secuencia ordenada de fases, desde
el reconocimiento inicial hasta la consecución del objetivo.

\item[LSASS] \textit{Local Security Authority Subsystem Service}. Proceso de Windows que
gestiona la autenticación y custodia credenciales en memoria; objetivo habitual de las
técnicas de robo de credenciales.

\item[MATCH\_RECOGNIZE] Construcción del estándar SQL para reconocer patrones de filas
ordenadas; Flink la emplea como lenguaje de procesamiento de eventos complejos.

\item[MulVAL] \textit{Multihost, multistage Vulnerability Analysis}. Analizador de
seguridad de red basado en lógica que emplea Datalog para derivar rutas de ataque
factibles a partir de la configuración y las vulnerabilidades de una red.

\item[OCSF] \textit{Open Cybersecurity Schema Framework}. Esquema abierto y agnóstico de
proveedor para la normalización de eventos de seguridad.

\item[Sigma] Formato abierto y agnóstico para expresar reglas de detección sobre registros
de eventos, convertible a los lenguajes de consulta de distintos SIEM.

\item[Oráculo] Versión de referencia de un patrón, calculada por lotes sobre el corpus completo,
contra la que se valida la versión en flujo.

\item[SIEM] \textit{Security Information and Event Management}. Plataforma de recolección,
correlación y análisis centralizado de eventos de seguridad.

\item[SOC] \textit{Security Operations Centre}. Centro de operaciones de seguridad: equipo
responsable de la monitorización, detección y respuesta ante incidentes.

\item[STIX] \textit{Structured Threat Information eXpression}. Lenguaje estandarizado por
OASIS para representar e intercambiar inteligencia de amenazas de forma estructurada.

\item[Sysmon] \textit{System Monitor}. Herramienta de Microsoft que amplía la auditoría
nativa de Windows con telemetría detallada de procesos, red, ficheros y registro.

\item[Topic] Canal con nombre del bus de eventos en el que los productores publican mensajes y
del que los consumidores los leen.

\item[TTP] \textit{Tactics, Techniques and Procedures}. Descripción del comportamiento de un
adversario en tres niveles de abstracción: el objetivo perseguido, el método empleado y su
implementación concreta.

\item[Ventana de sesión] Ventana temporal que agrupa eventos separados entre sí por menos de un
hueco dado y se cierra cuando ese hueco se supera.

\item[Watermark] Mecanismo de los sistemas de procesamiento en \textit{streaming} que estima
el progreso del tiempo de evento, permitiendo decidir cuándo una ventana temporal puede
considerarse completa.

\end{description}

\clearpage

\pagenumbering{arabic}

%%%%% CUERPO DEL DOCUMENTO %%%%%%%

\chapter{Introducción}
\label{cap:introduccion}

\section{Contexto y motivación}

La detección de intrusiones en entornos corporativos ha evolucionado en las últimas dos
décadas desde el análisis de firmas sobre tráfico de red hasta plataformas que agregan
telemetría de múltiples fuentes: cortafuegos, sistemas de detección de red, agentes de
\textit{endpoint}, servicios en la nube y registros de aplicación. Esa evolución ha
resuelto un problema de \textit{visibilidad}, pero ha creado otro de \textit{escala}.

Un centro de operaciones de seguridad (SOC) de tamaño medio recibe diariamente un volumen
de alertas que excede con claridad su capacidad de análisis. La consecuencia directa es la
llamada \textbf{fatiga de alertas}: cuando la proporción de avisos irrelevantes es
suficientemente alta, el analista deja de tratarlos individualmente y adopta estrategias de
descarte que, inevitablemente, acaban filtrando también señales legítimas. El problema no
es únicamente de cantidad, sino de \textbf{ausencia de contexto}. Una alerta que informa de
que un proceso ha accedido a la memoria de otro no dice, por sí sola, si ese comportamiento
corresponde a una técnica adversaria documentada, si ha sido observada previamente en
campañas de un actor conocido, ni qué activos de la organización quedan comprometidos si la
sospecha se confirma.

A esta situación se añade una segunda dificultad, de naturaleza estructural. Los ataques
dirigidos no se manifiestan como eventos aislados sino como \textbf{secuencias de etapas}
encadenadas en el tiempo. Un adversario que busca credenciales privilegiadas ejecuta primero
un reconocimiento del entorno, obtiene después acceso a un equipo, carga las herramientas
necesarias, extrae la información y, con frecuencia, borra sus huellas. Cada una de esas
acciones, considerada por separado, puede resultar indistinguible de la actividad
administrativa legítima. Es la \textit{secuencia}, y no cada evento individual, la que
constituye la evidencia sólida del ataque. Los sistemas de detección que evalúan eventos de
forma independiente son, por construcción, incapaces de reconocer esa estructura.

El panorama de amenazas actual agrava ambos problemas. Los informes de referencia del sector
señalan un crecimiento sostenido de los ataques asistidos por inteligencia artificial y una
consolidación de MITRE ATT\&CK \cite{mitre_attack} como lenguaje común para describir el
comportamiento adversario \cite{enisa2025}. Existe por tanto una necesidad clara de
\textbf{pasar de la alerta aislada al incidente correlacionado, contextualizado y
priorizado}, que es precisamente el ámbito de la denominada \textit{conciencia
cibersituacional}.

\section{Planteamiento del problema}

El problema que aborda este trabajo puede formularse en tres preguntas encadenadas:

\begin{enumerate}
    \item ¿Cómo detectar automáticamente que un conjunto de eventos dispersos, procedentes
    de fuentes heterogéneas, constituye en realidad una \textbf{secuencia de ataque
    coherente} y no actividad legítima?
    \item Una vez detectada, ¿cómo determinar el \textbf{alcance real del riesgo} que esa
    secuencia supone para la organización, en función de los activos que puede llegar a
    comprometer?
    \item ¿Cómo presentar el resultado de forma que el analista pueda \textbf{confiar en la
    priorización y comprender su fundamento}, en lugar de recibir una puntuación opaca?
\end{enumerate}

Las tres preguntas comparten un requisito transversal: la respuesta debe producirse en
tiempo útil sobre un flujo continuo de datos, no como resultado de un análisis
\textit{a posteriori}.

\section{Objetivos}

\subsection{Objetivo general}

Diseñar, implementar y validar una plataforma de datos que detecte ataques multi-etapa en
tiempo real y priorice los incidentes resultantes de forma explicable, integrando
inteligencia de amenazas y razonamiento sobre rutas de ataque para dar soporte a la
conciencia cibersituacional.

\subsection{Objetivos específicos}

\begin{enumerate}
    \item \textbf{Ingerir y normalizar} telemetría heterogénea de red y de sistema operativo
    a un esquema canónico común, sobre una arquitectura de procesamiento en
    \textit{streaming}.

    \item \textbf{Detectar secuencias de ataque multi-etapa} mediante procesamiento de
    eventos complejos, alineando las detecciones con las tácticas y técnicas de MITRE
    ATT\&CK.

    \item \textbf{Construir y mantener un grafo de ataque} que integre activos, topología,
    vulnerabilidades e inteligencia de amenazas estructurada.

    \item \textbf{Aplicar razonamiento lógico} sobre dicho grafo para derivar las rutas de
    ataque factibles hacia los activos críticos de la organización.

    \item \textbf{Priorizar los incidentes de forma explicable}, combinando la confianza de
    la detección, la severidad de la etapa alcanzada, la criticidad del activo afectado y su
    proximidad en una ruta de ataque derivada.

    \item \textbf{Validar cuantitativamente} el sistema sobre conjuntos de datos etiquetados
    con MITRE ATT\&CK, y proporcionar un despliegue reproducible.
\end{enumerate}

\section{Metodología de trabajo}

El desarrollo sigue un enfoque \textbf{incremental por fases}, en el que cada fase produce
un sistema funcional y verificable antes de acometer la siguiente. Esta decisión responde a
dos consideraciones prácticas: por un lado, garantiza que en cualquier momento del proyecto
exista una versión demostrable; por otro, permite validar las hipótesis técnicas de mayor
riesgo de forma temprana, antes de invertir esfuerzo sobre una arquitectura que pudiera
resultar inviable.

Cada componente se valida sobre datos reales de ataques documentados desde el momento de su
construcción, en lugar de posponer la verificación a una fase final de pruebas. Este
principio ha demostrado su utilidad durante el desarrollo: varios de los problemas más
relevantes descritos en el Capítulo~\ref{cap:desarrollo} son errores silenciosos que no
producían fallo alguno y que solo se detectaron al inspeccionar la coherencia de los
resultados intermedios.

Todas las decisiones de arquitectura relevantes se han registrado de forma continua junto
con su justificación y las alternativas descartadas, práctica que constituye la base del
apartado de decisiones de diseño de este documento.

\section{Estructura del documento}

El resto de la memoria se organiza como sigue:

\begin{itemize}
    \item El \textbf{Capítulo~\ref{cap:estadoarte}} revisa el estado del arte: los modelos y
    estándares sobre los que se apoya el trabajo (ATT\&CK, OCSF, STIX, Sigma), las técnicas
    de detección de secuencias, los grafos de ataque y el razonamiento lógico aplicado a
    seguridad, y los trabajos relacionados, identificando el hueco que este trabajo pretende
    cubrir.

    \item El \textbf{Capítulo~\ref{cap:desarrollo}} describe el análisis de requisitos, la
    arquitectura propuesta, las decisiones de diseño adoptadas y la implementación de cada
    uno de los componentes del sistema.

    \item El \textbf{Capítulo~\ref{cap:resultados}} presenta el protocolo de evaluación, los
    experimentos realizados y el análisis de los resultados obtenidos frente a las líneas
    base establecidas.

    \item El \textbf{Capítulo~\ref{cap:conclusiones}} recoge las conclusiones del trabajo y
    las líneas de continuación identificadas.

\end{itemize}

Los anexos recogen el material complementario que no cuenta para la extensión del cuerpo: los
aspectos éticos, sociales, económicos y ambientales (Anexo~\ref{cap:aspectos}), el presupuesto
(Anexo~\ref{cap:presupuesto}), la especificación del esquema de campaña
(Anexo~\ref{anx:esquema}), la implementación de los patrones en Flink SQL
(Anexo~\ref{anx:patrones}), el entorno y la reproducción de los experimentos
(Anexo~\ref{anx:reproducibilidad}) y las licencias de los componentes empleados
(Anexo~\ref{anx:licencias}).


\chapter{Estado del arte}
\label{cap:estadoarte}

Este capítulo revisa los fundamentos sobre los que se asienta el trabajo. Se organiza
siguiendo el recorrido que sigue un dato dentro del sistema propuesto: primero el contexto
operativo y el problema que se pretende resolver (Sección~\ref{sec:soc}); después los
modelos y estándares que permiten hablar un lenguaje común
(Secciones~\ref{sec:attack}--\ref{sec:sigma}); a continuación las técnicas de detección de
secuencias (Sección~\ref{sec:cep}), de razonamiento sobre grafos de ataque
(Sección~\ref{sec:grafos}) y de inteligencia de amenazas (Sección~\ref{sec:cti}); y
finalmente el marco integrador de la conciencia cibersituacional
(Sección~\ref{sec:conciencia}) y los trabajos relacionados
(Sección~\ref{sec:relacionados}).

\section{Operaciones de seguridad y fatiga de alertas}
\label{sec:soc}

Un centro de operaciones de seguridad articula su trabajo en torno a un ciclo de
monitorización, detección, análisis y respuesta. La pieza tecnológica central de ese ciclo
ha sido tradicionalmente el SIEM (\textit{Security Information and Event Management}), cuya
función es centralizar los registros procedentes de fuentes heterogéneas, normalizarlos,
aplicar reglas de correlación y presentar alertas al analista.

El modelo presenta dos limitaciones bien conocidas en la práctica operativa. La primera es
cuantitativa: el número de alertas generadas crece con la superficie monitorizada, mientras
que la capacidad de análisis permanece esencialmente constante. Cuando la proporción de
falsos positivos es elevada, el analista desarrolla estrategias de descarte que degradan la
eficacia global del sistema, fenómeno documentado como \textit{fatiga de alertas}.

La segunda limitación es cualitativa y, a efectos de este trabajo, más determinante: la
alerta se emite \textbf{desprovista de contexto}. El sistema informa de que se ha cumplido
una condición sobre un evento, pero no de qué significa ese evento en términos de
comportamiento adversario, ni de qué consecuencias tiene para los activos de la
organización. La consecuencia es que el trabajo de contextualización —consultar qué técnica
corresponde, si hay inteligencia asociada, qué máquinas están implicadas y qué podría
alcanzar el atacante desde ahí— recae íntegramente sobre el analista, de forma manual y
repetida para cada alerta.

Estas dos limitaciones motivan directamente los objetivos del presente trabajo: reducir el
volumen mediante correlación y aportar el contexto de forma automática.

% TODO: ampliar con referencias específicas sobre alert fatigue (estudios cuantitativos
% del ratio de falsos positivos en SOC reales) tras la revisión bibliográfica de agosto.

\section{MITRE ATT\&CK como lenguaje común}
\label{sec:attack}

MITRE ATT\&CK \cite{mitre_attack} es una base de conocimiento pública que cataloga el
comportamiento adversario observado en intrusiones reales, organizado en una jerarquía de
tres niveles:

\begin{itemize}
    \item \textbf{Tácticas}: el objetivo que persigue el adversario en una fase concreta
    (por ejemplo, \textit{Credential Access}, identificada como \texttt{TA0006}).
    \item \textbf{Técnicas}: el método empleado para lograr ese objetivo (por ejemplo,
    \textit{OS Credential Dumping}, \texttt{T1003}).
    \item \textbf{Subtécnicas}: la variante concreta de la técnica (por ejemplo,
    \textit{LSASS Memory}, \texttt{T1003.001}).
\end{itemize}

La relevancia de ATT\&CK para este trabajo es doble. Por un lado, proporciona un
\textbf{vocabulario compartido} que permite que la telemetría, las reglas de detección, la
inteligencia de amenazas y los informes de incidente se refieran a los mismos conceptos con
los mismos identificadores; sin ese vocabulario, la integración de las distintas fuentes que
plantea la arquitectura propuesta sería inviable. Por otro, la organización en tácticas
proporciona un \textbf{orden implícito de progresión} del ataque que resulta directamente
aprovechable para el modelado de secuencias: un adversario que ha alcanzado la fase de
exfiltración se encuentra objetivamente en una situación más grave que uno que solo ha
completado el reconocimiento.

\section{Normalización de telemetría: el esquema OCSF}
\label{sec:ocsf}

El problema de la heterogeneidad de formatos es transversal a toda la industria: cada
fabricante emite sus registros con nombres de campo, estructuras y semánticas propias. La
consecuencia práctica es que la lógica de detección queda acoplada al formato de origen, de
modo que una regla escrita para una fuente no es reutilizable con otra.

El \textit{Open Cybersecurity Schema Framework} (OCSF) \cite{ocsf} aborda este problema
mediante un esquema abierto y agnóstico de proveedor. Su modelo organiza los eventos en
\textbf{categorías} (actividad de sistema, hallazgos, gestión de identidades y accesos, red,
descubrimiento y actividad de aplicación), cada una de las cuales agrupa un conjunto de
\textbf{clases de evento} identificadas por un \texttt{class\_uid}. Así, por ejemplo, la
clase \textit{Process Activity} (\texttt{1007}) representa la creación, terminación, apertura
o inyección de procesos, con independencia de si el dato procede de Sysmon en Windows, de
\textit{auditd} en Linux o de un agente comercial.

Cada evento incorpora además un \texttt{activity\_id} que precisa la acción concreta dentro
de la clase, y un \texttt{type\_uid} derivado de ambos según la relación
$type\_uid = class\_uid \times 100 + activity\_id$. El esquema define asimismo un conjunto de
campos obligatorios comunes a todos los eventos y un campo \texttt{unmapped} destinado a
preservar la información de origen que no encuentra correspondencia en el modelo, lo que
evita la pérdida de datos durante la normalización.

La adopción de OCSF en este trabajo responde a una decisión de diseño explícita: desacoplar
la lógica de detección del formato de las fuentes, de modo que incorporar una fuente nueva
requiera únicamente escribir su mapeador correspondiente, sin modificar las reglas ni los
patrones de correlación situados aguas abajo.

\section{Reglas de detección: el formato Sigma}
\label{sec:sigma}

Sigma \cite{sigma} desempeña para las reglas de detección un papel análogo al que OCSF
desempeña para los datos: proporciona un formato abierto y agnóstico en el que expresar la
lógica de detección una sola vez, para después traducirla al lenguaje de consulta de cada
plataforma concreta mediante los denominados \textit{backends}.

Una regla Sigma se compone esencialmente de un bloque \texttt{logsource}, que declara sobre
qué tipo de registro aplica, y un bloque \texttt{detection}, que expresa las condiciones que
deben cumplirse. El repositorio público de la comunidad mantiene un catálogo extenso de
reglas, la mayoría de ellas \textbf{etiquetadas con la técnica ATT\&CK correspondiente}, lo
que las hace directamente aprovechables para el propósito de este trabajo.

Conviene señalar una dificultad de integración que resulta relevante para el desarrollo: las
reglas Sigma están escritas mayoritariamente contra los nombres de campo nativos de Windows
y Sysmon, mientras que la arquitectura propuesta opera sobre eventos ya normalizados a OCSF.
Su aplicación exige por tanto construir una tabla de equivalencia entre ambos modelos, que
es funcionalmente lo que implementa un \textit{backend} de Sigma para un SIEM comercial.

\section{Procesamiento de eventos complejos y detección multi-etapa}
\label{sec:cep}

El procesamiento de eventos complejos (CEP) es el paradigma que permite expresar y detectar
\textbf{patrones formados por secuencias de eventos} relacionados en el tiempo, en lugar de
evaluar cada evento de forma independiente. Un motor CEP mantiene estado sobre el flujo de
entrada y emite un resultado únicamente cuando la secuencia completa se confirma, lo que lo
convierte en la herramienta natural para el problema planteado en este trabajo.

Apache Flink \cite{flink} es un motor de procesamiento distribuido en \textit{streaming} que
proporciona dos vías para expresar este tipo de patrones. La primera es la
\textit{Pattern API} de FlinkCEP \cite{flinkcep}, que permite componer patrones mediante
cuantificadores, condiciones iterativas, ventanas temporales y condiciones de ausencia. La
segunda es la construcción \texttt{MATCH\_RECOGNIZE} de SQL, estandarizada en SQL:2016 bajo
el nombre de \textit{row pattern recognition}, que expresa la misma clase de patrones
mediante una sintaxis declarativa con las cláusulas \texttt{PATTERN}, \texttt{DEFINE},
\texttt{MEASURES} y \texttt{WITHIN}.

Dos conceptos del procesamiento en \textit{streaming} resultan críticos para la corrección
de este tipo de detecciones. El primero es la distinción entre \textbf{tiempo de evento}
(\textit{event time}), correspondiente al instante en que la acción ocurrió realmente, y
\textbf{tiempo de procesamiento} (\textit{processing time}), correspondiente al instante en
que el sistema la procesa; la detección de secuencias solo es correcta si se apoya en el
primero. El segundo son las \textbf{marcas de agua} (\textit{watermarks}), el mecanismo
mediante el cual el motor estima el progreso del tiempo de evento y decide cuándo una
ventana puede considerarse cerrada. Como se documenta en el Capítulo~\ref{cap:desarrollo},
la gestión incorrecta de este segundo mecanismo impide por completo la emisión de
resultados, sin producir error alguno que lo delate.

\section{Grafos de ataque y razonamiento lógico}
\label{sec:grafos}

Un \textbf{grafo de ataque} representa las formas en que un adversario puede progresar a
través de una red comprometiendo sucesivamente elementos de la misma, a partir de la
topología, las vulnerabilidades presentes y las políticas de acceso vigentes. Su utilidad
reside en que traslada el análisis desde el plano del evento individual al plano del
\textit{riesgo estructural}: no solo qué ha ocurrido, sino qué puede llegar a ocurrir como
consecuencia.

MulVAL \cite{mulval} constituye la referencia canónica en este ámbito. Su aportación
fundamental es adoptar \textbf{Datalog} como lenguaje de modelado: la configuración de la
red, las vulnerabilidades y las políticas se expresan como \textit{hechos}, mientras que la
forma en que un atacante puede aprovecharlas se expresa mediante \textit{reglas de
interacción}. El motor de inferencia deriva entonces automáticamente el conjunto de estados
alcanzables por el adversario, y por tanto las rutas de ataque factibles hacia los activos
críticos.

Este enfoque presenta tres propiedades especialmente valiosas para el objetivo de este
trabajo. En primer lugar, es \textbf{declarativo}: las reglas describen el conocimiento
sobre cómo progresa un ataque, no el procedimiento para calcularlo. En segundo lugar, es
\textbf{trazable}: cada conclusión derivada puede justificarse indicando la regla aplicada y
los hechos que la sustentan, lo que constituye la base natural de la explicabilidad
perseguida. Y en tercer lugar, es \textbf{extensible}: incorporar un nuevo mecanismo de
progresión adversaria consiste en añadir una regla, sin alterar el resto del sistema.

De forma complementaria, el modelado del grafo mediante bases de datos orientadas a grafos
de propiedades permite aplicar sobre él técnicas de \textbf{analítica de grafos}: medidas de
centralidad para identificar los nodos que actúan como puntos de paso críticos de la red,
algoritmos de caminos mínimos ponderados para estimar la ruta más probable, o detección de
comunidades para reconocer la segmentación efectiva. Esta combinación —razonamiento lógico
para derivar lo que es \textit{posible}, analítica de grafos para estimar lo que es
\textit{probable}— es la que adopta el presente trabajo.

% TODO: añadir referencias a extensiones recientes de MulVAL y a trabajos sobre generación
% automática de grafos de ataque a partir de informes CTI (revisión de agosto).

\section{Inteligencia de amenazas: STIX, plataformas y extracción automática}
\label{sec:cti}

La inteligencia de amenazas (CTI) proporciona el contexto que permite pasar de «se ha
observado este comportamiento» a «este comportamiento corresponde a una técnica empleada por
un actor conocido en campañas documentadas». Su representación estructurada se ha
estandarizado en torno a \textbf{STIX} \cite{stix}, un lenguaje que modela entidades como
actores, campañas, malware, patrones de ataque e indicadores, junto con las relaciones entre
ellas. Esa estructura de entidades y relaciones es, conceptualmente, un grafo de
conocimiento.

En el plano de las herramientas, OpenCTI \cite{opencti} constituye la plataforma de
referencia de código abierto para la gestión de este tipo de inteligencia. Su modelo de
datos es nativo en STIX~2.1 y dispone de conectores que permiten importar automáticamente la
base de conocimiento de MITRE ATT\&CK y sincronizar fuentes públicas de indicadores.

Un problema abierto y activo en este ámbito es la \textbf{extracción automática de TTP a
partir de informes en lenguaje natural}. La mayor parte de la inteligencia pública se publica
en forma de informes narrativos, no estructurados, lo que obliga a un trabajo manual de
lectura y anotación para incorporarla a una plataforma. En esta línea, TRAM
\cite{tram} aborda específicamente el mapeo automático de texto de informes a técnicas de
ATT\&CK mediante un modelo de lenguaje del dominio afinado para la tarea, lo que permite
cerrar el ciclo desde el informe publicado hasta el grafo consultable.

\section{Aprendizaje automático aplicado a la detección}
\label{sec:ml}

Los enfoques descritos hasta ahora —reglas de detección, patrones temporales, razonamiento
lógico— comparten una característica esencial: son \textbf{simbólicos}. Alguien codifica
explícitamente el conocimiento, y el sistema aplica exactamente lo codificado. Esto aporta
precisión y, sobre todo, explicabilidad: siempre es posible indicar qué condición se cumplió.

Su límite es igualmente evidente: \textbf{solo detectan lo que alguien ha descrito}. Una
técnica sin regla escrita pasa inadvertida, por llamativa que resulte su ejecución.

\subsection{Detección de anomalías}

La detección de anomalías invierte el planteamiento. En lugar de describir el comportamiento
malicioso, se modela el comportamiento \textit{normal} y se señala lo que se desvía de él.
La ventaja es que no requiere conocer la amenaza de antemano; el inconveniente es que
«distinto» no equivale a «malicioso», por lo que la tasa de falsos positivos es
estructuralmente alta.

Entre los métodos no supervisados, \textbf{Isolation Forest} \cite{isolationforest} resulta
particularmente adecuado para este dominio. A diferencia de las técnicas basadas en
distancias o densidades, aísla las observaciones mediante particiones aleatorias del espacio
de características, partiendo de la premisa de que las instancias anómalas requieren menos
particiones para quedar aisladas. Ello le confiere complejidad lineal y buen comportamiento
en espacios de dimensión elevada, además de permitir declarar explícitamente la proporción
esperada de anomalías, lo que resulta útil cuando el conjunto de entrenamiento no está
completamente libre de actividad maliciosa.

\subsection{El problema del desbalanceo}

La detección de intrusiones presenta un desbalanceo extremo por naturaleza: la actividad
maliciosa constituye una fracción mínima del total. Un clasificador que prediga siempre
«normal» alcanzaría una exactitud superior al 99\,\% siendo completamente inútil, lo que
obliga a emplear métricas robustas frente al desbalanceo —precisión, exhaustividad, área
bajo la curva precisión-exhaustividad— en lugar de la exactitud global.

Entre las técnicas de compensación, \textbf{SMOTE} \cite{smote} genera muestras sintéticas
de la clase minoritaria interpolando entre vecinos próximos, en lugar de replicar las
existentes. Su eficacia depende de la dimensionalidad: en espacios de muchas
características, la interpolación puede generar muestras poco representativas.

\subsection{Aprendizaje sobre grafos}

Cuando el objeto de análisis es un grafo, las técnicas de aprendizaje sobre grafos permiten
obtener representaciones vectoriales de los nodos que capturan su posición estructural. Sobre
esas representaciones puede plantearse la \textbf{predicción de enlaces}: estimar qué
relaciones son plausibles aunque no estén declaradas explícitamente.

Trasladado al grafo de ataque descrito en la Sección~\ref{sec:grafos}, esto permite sugerir
rutas de progresión que las reglas de interacción no contemplan porque nadie las formuló.

\subsection{La tensión entre generalización y explicabilidad}

Los enfoques simbólico y estadístico presentan fortalezas complementarias y debilidades
opuestas:

\begin{center}
\begin{tabular}{@{}lll@{}}
\toprule
 & \textbf{Simbólico} & \textbf{Aprendizaje automático} \\
\midrule
Cobertura & Solo lo codificado & Generaliza a lo no visto \\
Explicabilidad & Total y trazable & Limitada u opaca \\
Falsos positivos & Controlables por diseño & Estructuralmente altos \\
Mantenimiento & Requiere codificación manual & Reentrenamiento \\
\bottomrule
\end{tabular}
\end{center}

Esta complementariedad motiva los enfoques \textbf{neurosimbólicos}, que combinan ambos
paradigmas. Aplicado al problema que ocupa a este trabajo, el planteamiento consiste en que
el componente estadístico \textit{proponga} —hipótesis, relaciones plausibles, comportamiento
inusual— y el componente simbólico \textit{verifique y justifique}. De este modo se amplía la
cobertura sin renunciar a la capacidad de explicar las conclusiones, que es un requisito
explícito de este trabajo.

% TODO: completar con referencias específicas de IA neurosimbólica y de trabajos que la
% apliquen a seguridad (revisión bibliográfica pendiente).

\section{Conciencia cibersituacional}
\label{sec:conciencia}

Los elementos anteriores —detección de secuencias, grafos de ataque, inteligencia de
amenazas— encuentran su marco integrador en el concepto de \textbf{conciencia
cibersituacional}, entendido como la capacidad de percibir los elementos relevantes del
entorno, comprender su significado en conjunto y proyectar su evolución futura.

Trasladado al dominio que ocupa a este trabajo, ese esquema de tres niveles se corresponde
de forma directa con los tres pilares de la arquitectura propuesta:

\begin{itemize}
    \item La \textbf{percepción} corresponde a la detección de secuencias multi-etapa: saber
    qué está ocurriendo realmente, más allá de eventos sueltos.
    \item La \textbf{comprensión} corresponde al grafo de ataque razonado: entender qué
    significa lo detectado en términos de riesgo estructural para la organización.
    \item La \textbf{proyección} corresponde a la priorización explicable y a las rutas de
    ataque derivadas: anticipar hacia dónde puede evolucionar la situación y decidir en
    consecuencia.
\end{itemize}

Estrechamente ligado a este concepto se encuentra el de \textbf{gestión dinámica de
riesgos}, que propone sustituir la evaluación estática y periódica del riesgo por una
estimación continua que se actualiza con la evidencia observada en tiempo real. La
priorización que plantea este trabajo, al combinar la detección con la criticidad de los
activos y la proximidad en una ruta de ataque, se inscribe en esa línea.

\section{Trabajos relacionados y hueco identificado}
\label{sec:relacionados}

Cada uno de los componentes descritos cuenta con soluciones maduras consideradas de forma
aislada. Existen SIEM capaces de aplicar reglas sobre grandes volúmenes de telemetría;
existen motores CEP de propósito general con garantías de procesamiento sólidas; existen
generadores de grafos de ataque basados en lógica; y existen plataformas de inteligencia de
amenazas con modelos de datos estandarizados.

La observación que motiva este trabajo es que dichos componentes operan habitualmente
\textbf{desacoplados entre sí}. La detección no consulta el grafo de ataque; el grafo de
ataque se construye a partir de configuración estática y no incorpora la evidencia observada
en tiempo real; y la priorización de alertas, cuando existe, se basa en una severidad
declarada por la regla que la ha disparado, sin considerar el contexto del activo afectado
ni su exposición.

El hueco que este trabajo pretende cubrir es, por tanto, de \textbf{integración}: construir
una capa de correlación que enlace la detección de secuencias en tiempo real con el
razonamiento estructural sobre rutas de ataque, y que produzca a partir de esa unión una
priorización de incidentes justificada y comprensible. Las contribuciones concretas que se
persiguen son tres:

\begin{enumerate}
    \item Un \textbf{pipeline reproducible} que normaliza telemetría heterogénea a un
    esquema canónico y detecta sobre ella secuencias multi-etapa alineadas con ATT\&CK.
    \item Un \textbf{razonador de rutas de ataque} que incorpora, junto a la configuración y
    las vulnerabilidades, la evidencia efectivamente observada por la capa de detección.
    \item Un \textbf{modelo de priorización explicable} que descompone cada puntuación en
    factores interpretables y lo acompaña de una justificación legible.
    \item Una \textbf{integración de componentes de aprendizaje automático} en los puntos
    donde aportan capacidades que el enfoque simbólico no cubre, manteniendo la
    explicabilidad mediante verificación simbólica de sus propuestas.
\end{enumerate}

% TODO: completar esta sección con la comparativa explícita frente a trabajos concretos
% (tabla comparativa) una vez completada la revisión bibliográfica de agosto.


\chapter{Desarrollo}
\label{cap:desarrollo}

Este capítulo describe el análisis de requisitos, la arquitectura propuesta, las decisiones
de diseño adoptadas y la implementación del sistema. La exposición sigue el mismo orden
incremental en que se ha llevado a cabo el desarrollo, de modo que cada apartado se apoya en
las capacidades construidas en el anterior.

\section{Análisis de requisitos}

\subsection{Requisitos funcionales}

\begin{description}
    \item[RF1] El sistema debe ingerir telemetría de sistema operativo y de red procedente de
    fuentes heterogéneas, sobre un flujo continuo de datos.
    \item[RF2] El sistema debe normalizar los eventos ingeridos a un esquema canónico común,
    preservando la información de origen que no encuentre correspondencia en el modelo.
    \item[RF3] El sistema debe detectar técnicas adversarias individuales sobre el flujo
    normalizado, etiquetando cada detección con la técnica de MITRE ATT\&CK correspondiente.
    \item[RF4] El sistema debe detectar secuencias de ataque multi-etapa, esto es, patrones
    formados por varias detecciones relacionadas temporalmente.
    \item[RF5] El sistema debe mantener un grafo que integre activos, topología,
    vulnerabilidades e inteligencia de amenazas.
    \item[RF6] El sistema debe derivar rutas de ataque factibles hacia los activos críticos
    mediante razonamiento lógico sobre dicho grafo.
    \item[RF7] El sistema debe asignar a cada incidente una puntuación de riesgo acompañada
    de la descomposición de los factores que la determinan.
    \item[RF8] El sistema debe permitir la consulta analítica de los eventos almacenados con
    fines de investigación y evaluación.
    \item[RF9] El sistema debe mantener la trazabilidad de cada evento respecto de su origen,
    de forma que sea posible contrastar los resultados con la verdad conocida durante la
    evaluación.
\end{description}

\subsection{Requisitos no funcionales}

\begin{description}
    \item[RNF1] \textbf{Reproducibilidad}: el sistema completo debe poder desplegarse desde
    cero sobre una máquina de trabajo, sin intervención manual de configuración.
    \item[RNF2] \textbf{Corrección temporal}: el procesamiento debe apoyarse en el tiempo de
    ocurrencia real de los eventos, no en el de su procesamiento.
    \item[RNF3] \textbf{Ausencia de pérdida silenciosa}: todo evento que no pueda procesarse
    debe derivarse a una cola de rechazos con el motivo del fallo, nunca descartarse.
    \item[RNF4] \textbf{Observabilidad}: debe ser posible medir el estado y el rendimiento de
    cada componente del sistema en tiempo real.
    \item[RNF5] \textbf{Contención de recursos}: el sistema debe operar dentro de los límites
    de una estación de trabajo, sin requerir infraestructura dedicada.
    \item[RNF6] \textbf{Trazabilidad de decisiones}: las salidas del sistema deben ser
    explicables, evitando puntuaciones cuyo fundamento no pueda inspeccionarse.
\end{description}

\subsection{Restricciones}

El proyecto se apoya exclusivamente en herramientas de código abierto y conjuntos de datos
públicos. Toda la emulación de comportamiento adversario se realiza sobre datos previamente
capturados en entornos de laboratorio, sin ejecución sobre sistemas en producción. Los
recursos de cómputo se limitan a una estación de trabajo, lo que descarta por diseño
cualquier planteamiento que requiera el entrenamiento de modelos de gran tamaño.

\section{Visión general de la arquitectura}

La arquitectura se organiza en \textbf{dos planos de datos} que convergen en una
\textbf{capa de correlación y decisión}.

El \textbf{plano de telemetría} recorre el camino que va desde los sensores hasta las
detecciones: los eventos crudos se publican en un bus de eventos, se normalizan a un esquema
canónico, se evalúan contra reglas de detección y se correlacionan en busca de secuencias
multi-etapa. El \textbf{plano de inteligencia} construye en paralelo un grafo de conocimiento
a partir de fuentes de CTI y de la base de MITRE ATT\&CK. La capa de correlación cruza ambos:
enriquece las detecciones con el contexto del grafo, razona sobre las rutas de ataque
resultantes y emite la priorización final.

Sobre esa estructura se articulan los tres pilares que constituyen la aportación del trabajo:

\begin{description}
    \item[Pilar A --- Detección multi-etapa] Reconocimiento de secuencias de la
    \textit{kill chain} mediante procesamiento de eventos complejos, en lugar de alertas
    aisladas.
    \item[Pilar B --- Grafo de ataque razonado] Derivación de rutas de ataque factibles hacia
    los activos críticos mediante razonamiento lógico, con trazabilidad de cada paso.
    \item[Pilar C --- Priorización explicable] Asignación a cada incidente de una puntuación
    descompuesta en factores interpretables, acompañada de una justificación legible.
\end{description}

\begin{figure}[H]
\centering
\begin{tikzpicture}[
  caja/.style={draw=black!70,fill=black!6,rounded corners=2pt,minimum height=10mm,text width=19mm,align=center,font=\scriptsize,inner sep=2pt},
  hecho/.style={caja,fill=black!14},
  futuro/.style={caja,dashed,fill=white},
  topic/.style={font=\tiny\ttfamily},
  flecha/.style={-{Stealth[length=1.8mm]},thick},
  node distance=6mm and 6mm]
  \node[hecho] (datos) {Fuentes\\(Sysmon, Security, PowerShell)};
  \node[hecho,right=of datos] (norm) {Normalización\\OCSF};
  \node[hecho,right=of norm] (sigma) {Reglas Sigma\\y anomalías};
  \node[hecho,right=of sigma] (cep) {CEP en Flink\\(Pilar A)};
  \node[hecho,below=9mm of norm] (ch) {Almacén analítico\\(ClickHouse)};
  \node[futuro,below=9mm of cep] (grafo) {Grafo de ataque\\y Datalog (Pilar B)};
  \node[futuro,left=of grafo] (cti) {Inteligencia\\(ATT\&CK, STIX)};
  \node[futuro,below=9mm of grafo] (prio) {Priorización\\explicable (Pilar C)};
  \node[futuro,left=of prio] (api) {API y portal};
  \draw[flecha] (datos) -- node[above,topic]{raw} (norm);
  \draw[flecha] (norm) -- node[above,topic]{norm} (sigma);
  \draw[flecha] (sigma) -- node[above,topic]{det} (cep);
  \draw[flecha] (norm) -- (ch);
  \draw[flecha] (cep) -- node[right,topic]{campañas} (grafo);
  \draw[flecha] (cti) -- (grafo);
  \draw[flecha] (grafo) -- (prio);
  \draw[flecha] (prio) -- node[above,topic]{enr} (api);
  \begin{scope}[on background layer]
    \node[draw=black!40,dotted,rounded corners,fit=(datos)(norm)(sigma)(cep),inner sep=2mm,label={[font=\scriptsize]above:Plano de telemetría}] {};
  \end{scope}
\end{tikzpicture}
\caption{Arquitectura del sistema. Relleno gris: componentes construidos; borde discontinuo:
fases siguientes. Las etapas se comunican mediante \textit{topics} del bus de eventos.}
\label{fig:arquitectura}
\end{figure}

\subsection{Contrato de flujos de datos}

La comunicación entre etapas se realiza mediante \textit{topics} de un bus de eventos, con
una convención de nombres que refleja el grado de elaboración del dato: el prefijo
\texttt{raw} designa telemetría cruda, \texttt{norm} eventos ya normalizados, \texttt{det}
detecciones y \texttt{enr} incidentes enriquecidos. Los rechazos de cada etapa se derivan a
un \textit{topic} con prefijo \texttt{dlq}.

La Tabla~\ref{tab:topics} recoge el contrato establecido.

\begin{table}[H]
\centering
\begin{tabular}{|l|c|l|l|}
\hline
\textbf{Topic} & \textbf{Part.} & \textbf{Clave} & \textbf{Contenido} \\
\hline
\texttt{raw.host.windows} & 3 & \textit{hostname} & Telemetría cruda de sistema \\
\texttt{raw.net.zeek} & 3 & IP origen & Telemetría cruda de red \\
\texttt{norm.events} & 6 & \textit{hostname} & Eventos normalizados a OCSF \\
\texttt{det.alerts} & 3 & \textit{hostname} & Detecciones de reglas individuales \\
\texttt{det.campaigns} & 3 & \textit{hostname} & Secuencias multi-etapa detectadas \\
\texttt{enr.incidents} & 3 & \textit{id} incidente & Incidentes priorizados \\
\texttt{dlq.*} & 1 & --- & Rechazos, con el motivo del fallo \\
\hline
\end{tabular}
\caption{Contrato de \textit{topics} del sistema.}
\label{tab:topics}
\end{table}

La elección del \textbf{nombre de host como clave de particionado} responde a un requisito
funcional concreto: garantizar que todos los eventos de una misma máquina se procesen en
orden. La detección de secuencias carece de sentido si los eventos que la componen pueden
llegar desordenados, por lo que esta decisión, aparentemente menor, condiciona la viabilidad
del Pilar~A. Como se documenta en la Sección~\ref{sec:incidencia_watermark}, tiene además una
consecuencia no evidente sobre el procesamiento temporal.

\section{Decisiones de diseño}

Esta sección recoge las decisiones de arquitectura relevantes junto con su justificación y
las alternativas descartadas. Se documentan de forma explícita porque constituyen, en buena
medida, el contenido técnico del trabajo: en un proyecto de integración como este, el valor
reside tanto en las piezas elegidas como en los motivos de la elección.

\subsection{Representación del grafo: grafo de propiedades frente a RDF}

El planteamiento inicial contemplaba representar el grafo de conocimiento mediante RDF y
consultarlo con SPARQL, siguiendo la tradición de la web semántica. Esta opción se
\textbf{descartó} en favor de un grafo de propiedades consultado mediante Cypher, por tres
motivos. Primero, el grafo de propiedades es el modelo predominante en las herramientas de
análisis de rutas de ataque del ámbito de la seguridad ofensiva, lo que facilita la
comparación con el estado de la práctica. Segundo, la disponibilidad de bibliotecas de
analítica de grafos (centralidad, caminos mínimos, detección de comunidades) sobre este
modelo es considerablemente mayor. Y tercero, la capacidad de inferencia que aportaría una
ontología OWL queda cubierta —y superada para este dominio— por el razonador Datalog descrito
a continuación, que además es específico de seguridad.

\subsection{Razonamiento: Datalog frente a reglas ontológicas}

Para la derivación de rutas de ataque se adopta un motor Datalog siguiendo el enfoque de
MulVAL \cite{mulval}, en lugar de reglas de inferencia sobre una ontología. La razón es que
Datalog aplicado a análisis de seguridad de red es un enfoque canónico y consolidado, que
proporciona de forma natural la trazabilidad exigida por el requisito RNF6: cada conclusión
derivada puede justificarse indicando la regla aplicada y los hechos que la sustentan.

\subsection{Motor de procesamiento: Flink frente a micro-lotes}

Se adopta Apache Flink como motor de procesamiento en lugar de una solución basada en
micro-lotes. El motivo es que el Pilar~A requiere procesamiento con estado sobre tiempo de
evento y ventanas temporales precisas, capacidades que constituyen el núcleo del modelo de
Flink. Una arquitectura de micro-lotes obligaría a emular ese comportamiento con una
complejidad añadida considerable y garantías más débiles.

Dentro de Flink existían a su vez dos vías para expresar los patrones de secuencia: la
\textit{Pattern API} de FlinkCEP, disponible únicamente desde Java, y la construcción
\texttt{MATCH\_RECOGNIZE} de SQL, accesible también desde Python. Dado que el resto del
sistema se desarrolla en Python y que la segunda opción es un estándar SQL plenamente
expresivo para el tipo de patrones requeridos, se optó por esta última. La decisión se
validó experimentalmente antes de comprometerla, como se describe en la
Sección~\ref{sec:validacion_cep}.

\subsection{Componentes de infraestructura}

El resto de componentes se seleccionaron atendiendo a criterios de eficiencia operativa
sobre una estación de trabajo:

\begin{itemize}
    \item \textbf{Bus de eventos}: se emplea Redpanda, compatible con la API de Kafka pero
    implementado como binario único sin dependencia de JVM ni de un servicio de coordinación
    externo, lo que reduce sensiblemente el consumo de recursos y la complejidad operativa.
    \item \textbf{Almacén analítico}: se emplea ClickHouse, base de datos columnar orientada
    a analítica, por su rendimiento en consultas de agregación y filtrado sobre grandes
    volúmenes de eventos, que es el patrón de acceso característico de la investigación de
    incidentes.
    \item \textbf{Contenedores frente a máquinas virtuales}: todos los servicios se despliegan
    como contenedores sobre un único anfitrión. Dedicar una máquina virtual completa a cada
    servicio habría supuesto multiplicar el consumo de memoria por el sistema operativo
    invitado sin aportar aislamiento adicional relevante para este caso de uso.
\end{itemize}

\subsection{Papel del aprendizaje automático}

El aprendizaje automático se incorpora \textbf{únicamente allí donde aporta capacidades que
el enfoque simbólico no cubre}. Este criterio descarta expresamente los modelos que
duplicarían funciones ya resueltas: un clasificador que replicase el comportamiento de una
regla de detección añadiría opacidad sin ganar cobertura.

Bajo ese principio se identifican cuatro puntos de aplicación:

\begin{enumerate}
    \item \textbf{Detección de anomalías} sobre el flujo normalizado, que complementa a las
    reglas señalando comportamiento inusual sin regla asociada.
    \item \textbf{Extracción de técnicas} a partir de informes de amenazas en lenguaje
    natural, tarea inabordable mediante reglas.
    \item \textbf{Aprendizaje sobre el grafo de ataque}, para sugerir rutas de progresión no
    contempladas explícitamente por las reglas de interacción.
    \item \textbf{Ajuste del modelo de puntuación} a partir de datos, en lugar de pesos
    fijados manualmente.
\end{enumerate}

La restricción operativa es que el entrenamiento debe ser viable sobre la estación de
trabajo descrita, lo que acota el tamaño de los modelos empleados pero no impide el ajuste
fino de modelos del dominio ni el entrenamiento de modelos clásicos y de grafos.

En todos los casos se preserva la explicabilidad mediante el principio expuesto en la
Sección~\ref{sec:ml}: el componente estadístico propone y el simbólico verifica. Ninguna
salida del sistema se presenta al analista sin una justificación inspeccionable.

\section{Entorno de desarrollo}

El desarrollo se realiza sobre una estación de trabajo con procesador de 16 núcleos, 64~GB de
memoria y GPU dedicada, ejecutando el subsistema de Windows para Linux (WSL2) con Ubuntu como
sistema anfitrión de contenedores. Los límites de recursos asignados al subsistema
(44~GB de memoria y 24 hilos) se fijan de forma declarativa, reservando el resto para el
sistema operativo anfitrión.

Los servicios se orquestan mediante Docker Compose con \textbf{perfiles por fase}, lo que
permite levantar únicamente el subconjunto de componentes necesario para la tarea en curso.
Esta organización responde al requisito RNF5: el sistema completo excede los recursos
razonables de una estación de trabajo si todos los servicios se ejecutan simultáneamente,
mientras que por perfiles el consumo se mantiene contenido.

\begin{table}[H]
\centering
\begin{tabular}{|l|l|l|}
\hline
\textbf{Perfil} & \textbf{Servicios} & \textbf{Uso} \\
\hline
\texttt{core} & Bus de eventos y consola & Siempre activo \\
\texttt{stream} & Motor de procesamiento & Desarrollo de detección \\
\texttt{store} & Almacén analítico & Consulta y evaluación \\
\texttt{graph} & Base de datos de grafos & Trabajo sobre el Pilar B \\
\texttt{cti} & Plataforma de inteligencia & Ingesta de CTI \\
\texttt{observability} & Métricas y visualización & Monitorización \\
\hline
\end{tabular}
\caption{Perfiles de despliegue definidos.}
\label{tab:perfiles}
\end{table}

\subsection{Incidencias del entorno}

La construcción del entorno presentó una cadena de incompatibilidades derivadas de emplear
versiones muy recientes del sistema operativo y del intérprete de Python, que se documentan
por su valor práctico:

\begin{itemize}
    \item La distribución de Linux empleada es posterior a las versiones para las que existen
    repositorios oficiales de algunos componentes, lo que obligó a fijar explícitamente el
    repositorio de la última versión con soporte a largo plazo disponible.
    \item La versión de Python del sistema (3.14) carece de paquetes precompilados para las
    bibliotecas de integración con el motor de procesamiento, cuyas dependencias no soportan
    aún esa versión. Se resolvió mediante un entorno virtual independiente con Python 3.12,
    manteniendo el entorno principal para el resto del código.
    \item La biblioteca de integración requiere una máquina virtual Java en el anfitrión, no
    solo en los contenedores, lo que no resulta evidente hasta que falla la primera
    invocación.
\end{itemize}

\section{Ingesta de datos}

\subsection{Selección de los conjuntos de datos}

La validación requiere telemetría de ataques reales con \textbf{verdad conocida}, esto es,
con etiquetado explícito de las técnicas presentes. Se emplean para ello los conjuntos de
datos abiertos de OTRF \cite{otrf}, que reúnen capturas de laboratorio de técnicas
adversarias con dos propiedades especialmente valiosas: cada conjunto está mapeado a las
técnicas de MITRE ATT\&CK correspondientes, y las capturas incluyen no solo los eventos del
ataque sino también la actividad benigna concurrente, lo que permite evaluar falsos
positivos sobre datos realistas.

El repositorio organiza los conjuntos en dos familias. Los \textbf{atómicos} corresponden a
una única técnica aislada y resultan adecuados para la validación de componentes; los
\textbf{compuestos} recogen cadenas de ataque completas de varias etapas, y constituyen el
material de evaluación del Pilar~A.

Para el desarrollo inicial se seleccionó el conjunto atómico correspondiente al volcado de
memoria del proceso LSASS mediante la biblioteca \texttt{comsvcs.dll}, identificado como
\texttt{SDWIN-201018195009} y etiquetado con la técnica \texttt{T1003.001} de la táctica
\texttt{TA0006}. La elección responde a que, pese a tratarse de un conjunto pequeño
(184~eventos), contiene una cadena de acciones representativa y dos fuentes de auditoría
independientes que registran el mismo comportamiento, lo que permite validar la correlación
cruzada.

\subsection{Reproducción de la telemetría}

La ingesta se implementa mediante un componente de reproducción que lee los conjuntos de
datos —almacenados en formato JSON Lines, un evento por línea— y los publica en el bus de
eventos. Su diseño incorpora cuatro decisiones relevantes:

\begin{enumerate}
    \item \textbf{Clave de particionado por host}, conforme al contrato de la
    Tabla~\ref{tab:topics}, para garantizar el orden por máquina.
    \item \textbf{Preservación del tiempo de evento original}: la marca temporal del evento
    se propaga como marca temporal del mensaje en el bus, lo que permite que las etapas
    posteriores operen sobre tiempo de evento y no sobre tiempo de procesamiento (RNF2).
    \item \textbf{Trazabilidad mediante cabeceras}: cada mensaje incorpora el identificador
    del conjunto de datos de origen, lo que satisface el requisito RF9 y permitirá contrastar
    las detecciones con la verdad conocida durante la evaluación.
    \item \textbf{Control de velocidad de reproducción}: el componente admite reproducir a
    máxima velocidad, útil para pruebas, o respetando el ritmo temporal original de la
    captura, útil para demostraciones realistas del comportamiento en tiempo real.
\end{enumerate}

La reproducción del conjunto seleccionado se completó con la totalidad de los 184~eventos
publicados sin pérdidas, a un rendimiento aproximado de 3.300~eventos por segundo.

\section{Normalización a OCSF}

\subsection{Selección del subconjunto del esquema}

OCSF define un número elevado de clases de evento, de las cuales solo un subconjunto resulta
pertinente para la telemetría considerada. La selección se realizó a partir de las clases
efectivamente presentes en los conjuntos de datos, y se organizó en tres oleadas de
implementación para permitir que el sistema fuese funcional desde la primera.

\begin{table}[H]
\centering
\begin{tabular}{|l|c|l|}
\hline
\textbf{Clase OCSF} & \textbf{\texttt{class\_uid}} & \textbf{Fuentes que absorbe} \\
\hline
File System Activity & 1001 & Creación y escritura de ficheros \\
Module Activity & 1005 & Carga de bibliotecas dinámicas \\
Scheduled Job Activity & 1006 & Tareas programadas y servicios \\
Process Activity & 1007 & Creación, terminación y acceso a procesos \\
Event Log Activity & 1008 & Manipulación del registro de auditoría \\
Script Activity & 1009 & Ejecución de guiones \\
Account Change & 3001 & Alta y modificación de cuentas \\
Authentication & 3002 & Inicios de sesión \\
Authorize Session & 3003 & Asignación de privilegios especiales \\
Group Management & 3006 & Modificación de pertenencia a grupos \\
Network Activity & 4001 & Conexiones de red \\
DNS Activity & 4003 & Resolución de nombres \\
\hline
\end{tabular}
\caption{Subconjunto de clases OCSF adoptado.}
\label{tab:ocsf}
\end{table}

Merece una mención específica el tratamiento de los eventos de \textbf{registro de Windows}.
El esquema base de OCSF no define una clase específica para esta categoría, por lo que se
optó por asignarlos a la clase más próxima conservando la totalidad del detalle original en
el campo \texttt{unmapped}. La decisión se documenta como tal, dado que una futura extensión
del esquema podría cubrir este caso de forma más precisa.

\subsection{Arquitectura de los mapeadores}

La normalización se estructura en tres niveles. Un módulo común construye el esqueleto OCSF
compartido por todos los eventos: identificadores de clase y actividad, marca temporal,
metadatos, dispositivo de origen y los dos campos de extensión propios descritos más
adelante. Sobre ese esqueleto, un conjunto de \textbf{mapeadores específicos} —uno por
familia de origen y tipo de evento— añade los campos particulares de cada clase. Finalmente,
un \textbf{despachador} enruta cada evento entrante a su mapeador correspondiente en función
del par formado por el canal de origen y el identificador de evento.

Cuando el despachador no encuentra un mapeador para un evento, lanza una excepción específica
que la etapa de procesamiento traduce en la derivación del evento a la cola de rechazos,
junto con el motivo. Este mecanismo materializa el requisito RNF3: no existe descarte
silencioso.

\subsection{Extensiones propias del esquema}

El modelo incorpora dos extensiones sobre el esquema estándar:

\begin{itemize}
    \item El campo \texttt{unmapped}, previsto por el propio OCSF, conserva toda la
    información de origen sin correspondencia en el modelo canónico.
    \item Un bloque de extensión propio recoge la \textbf{trazabilidad} (identificador del
    conjunto de datos, canal y tipo de evento de origen) y un conjunto de \textbf{señales de
    detección precalculadas}.
\end{itemize}

La segunda extensión merece justificación, pues constituye una decisión de diseño con
consecuencias importantes aguas abajo. Durante el mapeo, el sistema evalúa un conjunto de
condiciones características de comportamiento adversario y las materializa como campos
booleanos: acceso al proceso que custodia las credenciales, acceso con permisos compatibles
con la lectura de memoria, carga de bibliotecas asociadas al volcado de procesos, escritura de
ficheros de volcado, borrado del registro de auditoría o inicio de sesión remoto.

El beneficio de precalcular estas señales en la normalización es doble. Por un lado, evita que
cada regla de detección deba reinterpretar los campos de origen, simplificando notablemente su
expresión. Por otro, y más relevante, constituye la \textbf{semilla de la explicabilidad}
perseguida en el Pilar~C: cuando el sistema priorice un incidente, podrá indicar exactamente
qué señales concretas contribuyeron a ello.

\subsection{Resultados de cobertura}

La cobertura de normalización —proporción de eventos que encuentran mapeador— se adoptó como
métrica de progreso de esta fase. La Tabla~\ref{tab:cobertura} recoge su evolución sobre el
conjunto de datos de referencia conforme se incorporaban mapeadores.

\begin{table}[H]
\centering
\begin{tabular}{|l|c|c|c|}
\hline
\textbf{Iteración} & \textbf{Normalizados} & \textbf{Rechazados} & \textbf{Cobertura} \\
\hline
Mapeadores iniciales (Sysmon 1/3/5/10/11) & 74 & 110 & 40{,}2\,\% \\
+ Módulos y registro (Sysmon 7/12/13/14) & 148 & 36 & 80{,}4\,\% \\
+ Canal Security (11 tipos de evento) & 184 & 0 & 100{,}0\,\% \\
\hline
\end{tabular}
\caption{Evolución de la cobertura de normalización sobre el conjunto \texttt{SDWIN-201018195009} (184 eventos).}
\label{tab:cobertura}
\end{table}

El proceso ilustra una propiedad deseable del diseño: la propia métrica indica qué construir a
continuación, pues los rechazos se agrupan por el par (canal, tipo de evento) que los provoca.
El resultado final son 12~tipos de evento OCSF distintos generados a partir de un único
conjunto de datos.

La equivalencia entre el flujo por lotes y el flujo en \textit{streaming} se verificó
ejecutando la misma lógica en ambos modos y contrastando los resultados, que coincidieron
exactamente.

\subsection{Correlación cruzada entre fuentes}

Un resultado relevante de esta fase es la confirmación independiente del comportamiento
adversario por parte de dos subsistemas de auditoría distintos. El acceso al proceso LSASS
queda registrado simultáneamente por la telemetría extendida de Sysmon y por la auditoría
nativa de Windows mediante los eventos de manipulación de descriptores de objeto.

Que dos mecanismos independientes registren el mismo hecho tiene un valor operativo directo:
constituye evidencia corroborada frente a un posible falso positivo, y ofrece resiliencia
ante la desactivación de una de las fuentes por parte del adversario.

Adicionalmente, la traza de llamadas preservada en el campo \texttt{unmapped} conserva la
secuencia completa de bibliotecas implicadas en la operación de volcado, lo que constituye la
evidencia forense que confirma la técnica concreta empleada.

\subsection{Incidencia: sincronización temporal entre fuentes}
\label{sec:incidencia_temporal}

Durante la validación se detectó un error de consecuencias potencialmente graves que no
producía fallo alguno. Los dos canales de telemetría considerados exponen campos temporales
distintos: la telemetría extendida incorpora un campo con la marca de tiempo en formato UTC
además del campo estándar, mientras que la auditoría nativa solo dispone de este último. La
implementación inicial priorizaba el campo aparentemente más preciso cuando estaba
disponible.

El resultado fue un desfase de dieciséis horas entre ambos canales que reordenaba
completamente la secuencia reconstruida: la etapa final del ataque —el borrado del registro de
auditoría— aparecía como la primera. El error no generaba ningún rechazo ni advertencia, y
solo se hizo visible al inspeccionar la coherencia lógica de la cadena de eventos.

La corrección consistió en apoyar la normalización en el campo temporal \textbf{común y
coherente} entre fuentes, en lugar del más específico. La lección resulta generalizable: en un
sistema que correlaciona múltiples fuentes, la homogeneidad del criterio temporal prevalece
sobre la precisión nominal de cada fuente considerada por separado.

Tras la corrección, la secuencia reconstruida presenta un entrelazado coherente de ambas
fuentes con resolución de milisegundos, condición necesaria para la detección de secuencias
del Pilar~A.

\section{Persistencia analítica}

El almacenamiento de los eventos normalizados persigue dos objetivos: soportar la
investigación interactiva sobre los datos y proporcionar el sustrato para la evaluación
cuantitativa del sistema.

\subsection{Diseño del esquema}

El diseño adopta una estrategia \textbf{híbrida}. Los campos que participan habitualmente en
filtros y agregaciones se promocionan a columnas dedicadas con tipos específicos, mientras que
el evento OCSF completo se conserva íntegro en una columna adicional. De este modo se obtiene
el rendimiento de una base de datos columnar sobre los accesos frecuentes sin renunciar a la
flexibilidad de disponer del dato completo cuando se requiere un campo no previsto.

Las decisiones de modelado más relevantes son las siguientes:

\begin{itemize}
    \item La \textbf{clave de ordenación} se establece sobre el par (host, tiempo), que
    corresponde al patrón de acceso dominante en la investigación de incidentes: qué ocurrió
    en una máquina concreta durante una ventana temporal determinada.
    \item El \textbf{particionado} se realiza por mes, lo que facilita la gestión y el purgado
    del almacén por rangos temporales.
    \item Las \textbf{señales de detección} precalculadas durante la normalización se
    materializan como columnas booleanas, lo que reduce las consultas de investigación sobre
    comportamiento adversario a filtros triviales.
    \item Se aplican \textbf{codificaciones específicas} por columna: diferencial más
    compresión para las marcas temporales, que presentan valores próximos entre sí, y
    diccionario para los campos de baja cardinalidad como nombres de host o de clase.
\end{itemize}

\subsection{Consultas de investigación}

La materialización de las señales como columnas permite reconstruir la cadena completa del
ataque mediante una única consulta declarativa que proyecta cada evento sobre la etapa
correspondiente y los ordena temporalmente. El resultado sobre el conjunto de referencia
recupera la secuencia esperada: preparación mediante carga de bibliotecas de volcado, acceso
al proceso de credenciales confirmado por ambas fuentes de auditoría, escritura del volcado a
disco y borrado del registro de auditoría.

Conviene señalar que esta reconstrucción se obtiene \textbf{antes de aplicar regla de
detección alguna}, únicamente a partir de las señales calculadas durante la normalización, lo
que valida el enfoque adoptado en esa fase.

\subsection{Incidencia: exposición del servicio de almacenamiento}

Se documenta una segunda incidencia por su carácter no evidente. El servicio de
almacenamiento, ejecutado en contenedor con el puerto correspondiente publicado, rechazaba
las conexiones procedentes del anfitrión pese a responder correctamente a las peticiones
originadas dentro del propio contenedor. El síntoma resultaba engañoso, ya que el error
devuelto sugería un problema de red o de redirección de puertos.

El diagnóstico se alcanzó comprobando que el servicio tampoco respondía a peticiones dirigidas
a su propia dirección de red: escuchaba exclusivamente en la interfaz de bucle local interna.
La solución consistió en declarar explícitamente la interfaz de escucha en la configuración
del servidor.

\section{Observabilidad}

La instrumentación del sistema responde al requisito RNF4 y persigue dos finalidades:
verificar durante el desarrollo que cada componente se comporta como se espera, y
proporcionar durante la evaluación las medidas de rendimiento que se reportan en el
Capítulo~\ref{cap:resultados}.

La recolección de métricas se realiza mediante un sistema de series temporales que consulta
periódicamente los puntos de exposición de cada componente. Se instrumentaron el bus de
eventos, ambos nodos del motor de procesamiento y el almacén analítico, activando en cada
caso el mecanismo de exposición correspondiente.

Sobre esa base se construyeron tres cuadros de mando, aprovisionados de forma
\textbf{declarativa} —definidos como ficheros de configuración versionados, no configurados
manualmente— para dar cumplimiento al requisito de reproducibilidad RNF1:

\begin{description}
    \item[Salud del \textit{pipeline}] Estado del bus de eventos, volumen por \textit{topic},
    recursos consumidos por el motor de procesamiento y actividad del almacén analítico.
    Verifica el correcto funcionamiento de la infraestructura.
    \item[Panel de seguridad] Distribución de eventos por clase OCSF, señales de detección
    activas, cronología del ataque reconstruida y procesos más frecuentes. Constituye la
    vista orientada al analista.
    \item[Cobertura y calidad de datos] Cobertura de normalización, distribución por canal y
    tipo de evento, conjuntos de datos cargados e indicadores de integridad temporal. Respalda
    las métricas reportadas en la evaluación.
\end{description}

\section{Validación experimental del enfoque de detección de secuencias}
\label{sec:validacion_cep}

Antes de comprometer el desarrollo del Pilar~A con la vía tecnológica seleccionada, se realizó
una validación experimental específica sobre datos reales. El objetivo era confirmar que la
construcción SQL elegida permitía efectivamente detectar secuencias sobre el flujo procedente
del bus de eventos, respetando el tiempo de evento original.

La prueba consistió en definir un patrón mínimo de dos eventos consecutivos particionado por
host, sobre la telemetría ya ingerida. El resultado inicial fue la ausencia total de
coincidencias, lo que motivó el diagnóstico descrito a continuación.

\subsection{Incidencia: bloqueo del avance temporal por particiones inactivas}
\label{sec:incidencia_watermark}

La causa resultó ser una consecuencia no evidente de la decisión de particionado adoptada. Al
emplear el nombre de host como clave y contener el conjunto de datos un único host, la
totalidad de los eventos se concentra en una sola partición, quedando las restantes vacías.

El motor de procesamiento calcula la marca de agua global como el mínimo de las marcas de agua
de todas sus fuentes. Las particiones que no reciben datos nunca emiten marca de agua alguna,
por lo que el mínimo permanece indefinido y \textbf{el tiempo de evento no avanza jamás}. La
consecuencia es que ninguna ventana temporal llega a considerarse cerrada y el motor no emite
resultado alguno, sin producir error que lo delate.

El diagnóstico se alcanzó por descarte, reduciendo el patrón a un único evento —que debería
haber emitido una coincidencia por cada evento de entrada— y comprobando que tampoco producía
salida, lo que descartaba un problema en la expresión del patrón y señalaba al mecanismo
temporal. La solución consiste en configurar un tiempo de inactividad tras el cual una fuente
sin datos se marca como inactiva y deja de bloquear el cálculo.

Resuelta la incidencia, el patrón de dos eventos emitió correctamente las coincidencias
esperadas, encadenadas y con las marcas temporales originales de la captura. Con ello quedó
validada la vía tecnológica y descartada la necesidad de recurrir a la alternativa prevista
como plan de contingencia.

Esta incidencia ilustra una interacción no trivial entre dos decisiones de diseño
independientes: el particionado por host, adoptado para garantizar el orden de eventos, tiene
como efecto colateral la inutilización del mecanismo de avance temporal cuando el número de
hosts es inferior al de particiones.

\section{Detección basada en reglas}
\label{sec:deteccion-reglas}

Sobre el flujo de eventos normalizados se construye la primera capa de detección, encargada
de identificar técnicas adversarias individuales. Se apoya en el catálogo público de reglas
Sigma descrito en la Sección~\ref{sec:sigma}, lo que permite aprovechar el trabajo de la
comunidad en lugar de redactar reglas propias.

\subsection{Traducción de las reglas al esquema canónico}

La integración presenta un obstáculo no evidente: las reglas del repositorio público están
escritas contra los nombres de campo \textbf{nativos} de Windows y Sysmon, mientras que el
sistema opera sobre eventos ya normalizados a OCSF. Aplicarlas directamente resulta
imposible, pues los campos que referencian no existen en el modelo canónico.

La biblioteca pySigma \cite{pysigma} prevé este escenario mediante el concepto de
\textit{pipeline de procesamiento}: una especificación declarativa que transforma las reglas
para adaptarlas a un modelo de datos concreto. Existe una implementación oficial para OCSF,
que se tomó como punto de partida.

Su análisis reveló, no obstante, \textbf{tres desviaciones} respecto de la semántica del
estándar que la hacían inaplicable al modelo de este trabajo:

\begin{enumerate}
    \item En los eventos de \textit{acceso a proceso}, el proceso origen se mapea al campo
    correspondiente al \textbf{proceso padre}. La correspondencia es incorrecta: en este tipo
    de evento el proceso origen no engendró al destino, sino que ejecuta una acción sobre él.
    OCSF define para ello el concepto de \textit{actor}, que es la correspondencia adecuada.
    \item Los eventos de \textit{carga de módulo} se clasifican como actividad de proceso con
    actividad genérica, cuando el esquema define una clase específica para ello
    (\textit{Module Activity}).
    \item Los accesos a proceso se etiquetan con la actividad genérica «otros» en lugar de la
    correspondiente a la apertura de un descriptor.
\end{enumerate}

La primera desviación produce detecciones incorrectas; las dos últimas impiden que las reglas
lleguen siquiera a evaluarse, dado que la condición sobre el identificador de tipo de evento
nunca se satisface.

Se optó por \textbf{implementar una pipeline propia} en lugar de extender la oficial. El
motivo es técnico: las condiciones que las pipelines añaden a las reglas se acumulan de forma
conjuntiva, por lo que una corrección superpuesta habría exigido simultáneamente dos valores
distintos del mismo campo, anulando toda posibilidad de coincidencia.

La pipeline resultante cubre nueve categorías de origen y establece la correspondencia entre
los campos nativos y el modelo OCSF del proyecto, dirigiendo al bloque de campos no mapeados
aquellos sin equivalente en el esquema. Su desarrollo constituye, en la práctica, la
implementación de un \textit{backend} de Sigma específico para este sistema.

\subsection{Evaluación de las reglas}

Para aplicar las reglas se plantearon dos vías. La primera consiste en \textbf{generar
consultas} en el lenguaje del almacén analítico y ejecutarlas periódicamente; es directa,
pero opera necesariamente en modo diferido. La segunda consiste en \textbf{evaluar la
estructura de la regla directamente contra cada evento} en memoria, sin lenguaje intermedio,
lo que permite el funcionamiento en tiempo real.

Se implementaron ambas: la primera como validación cruzada y la segunda como componente
operativo del sistema.

\subsubsection{Alcance de los modificadores soportados}

El formato Sigma admite un conjunto amplio de modificadores que alteran la comparación de
valores. Implementarlos todos constituiría un proyecto en sí mismo, por lo que se acotó el
alcance a partir de una medición sobre el repositorio completo de reglas.

\begin{table}[H]
\centering
\begin{tabular}{|l|r|r|}
\hline
\textbf{Modificador} & \textbf{Apariciones} & \textbf{Acumulado} \\
\hline
\texttt{contains}    & 3.362 & 45{,}1\,\% \\
\texttt{endswith}    & 2.494 & 78{,}5\,\% \\
\texttt{all}         &   854 & 90{,}0\,\% \\
\texttt{startswith}  &   431 & 95{,}8\,\% \\
\texttt{re}          &   155 & 97{,}9\,\% \\
\hline
Resto (\texttt{windash}, \texttt{cidr}, otros) & 158 & 100\,\% \\
\hline
\end{tabular}
\caption{Frecuencia de uso de los modificadores de Sigma en el repositorio público.}
\label{tab:modificadores}
\end{table}

Cinco modificadores cubren el 97,9\,\% de los usos reales. El evaluador los implementa
todos, junto con la lógica de composición de condiciones (conjunción, disyunción, negación
y cuantificación). Las construcciones no soportadas se detectan explícitamente y la regla
correspondiente se omite, en lugar de producir una evaluación silenciosamente incorrecta.

\subsubsection{Optimización}

La implementación inicial reconstruía el árbol de condiciones de cada regla en cada
evaluación. Sobre un conjunto de 2.093 reglas, ello suponía repetir más de 385.000 veces un
trabajo que puede realizarse una sola vez por regla.

Se introdujeron dos optimizaciones:

\begin{itemize}
    \item \textbf{Precompilación}: el árbol de condiciones se construye una única vez al
    cargar la regla.
    \item \textbf{Encaminamiento por tipo de evento}: cada regla declara la categoría de
    registro sobre la que aplica, lo que permite descartar de antemano aquellas que no pueden
    coincidir con el tipo del evento evaluado.
\end{itemize}

El efecto combinado se muestra en la Tabla~\ref{tab:optimizacion}. La comprobación esencial
es que el conjunto de detecciones producidas permaneció \textbf{idéntico}: se trata de una
optimización, no de una aproximación.

\begin{table}[H]
\centering
\begin{tabular}{|l|r|r|r|}
\hline
\textbf{Versión} & \textbf{Tiempo} & \textbf{Mejora} & \textbf{Detecciones} \\
\hline
Sin optimizar & 108{,}6 s & --- & 10 \\
Optimizada    & 1{,}49 s & 73$\times$ & 10 \\
\hline
\end{tabular}
\caption{Efecto de las optimizaciones sobre el conjunto de referencia (184 eventos, 2.093 reglas).}
\label{tab:optimizacion}
\end{table}

\subsection{Esquema de alerta}

Cada coincidencia produce una alerta que incorpora la identificación de la regla, su nivel de
severidad, las técnicas y tácticas de ATT\&CK asociadas, el dispositivo y usuario implicados,
la referencia al evento de origen y la trazabilidad al conjunto de datos.

El campo determinante es \texttt{matched\_fields}, que registra \textbf{qué condiciones
concretas dispararon la detección} con sus valores. Constituye la base de la explicabilidad
exigida por el requisito RNF6: la alerta no se limita a afirmar que algo es sospechoso, sino
que aporta la evidencia. A modo de ilustración, sobre el conjunto de referencia el sistema
recupera la línea de comandos completa empleada por el adversario, que coincide literalmente
con la documentada en los metadatos del conjunto de datos.

\subsection{Clasificación de las etiquetas ATT\&CK}

Durante la explotación de los resultados se detectó que las etiquetas del repositorio de
reglas no son homogéneas: bajo un mismo prefijo conviven identificadores de \textbf{técnicas},
de \textbf{software}, de \textbf{grupos de actores} y nombres de \textbf{táctica}, incluyendo
algunos no pertenecientes al estándar y acuñados por la propia comunidad. Una integración que
asuma uniformidad clasifica erróneamente estas etiquetas, contaminando el mapeo ATT\&CK.

El sistema las separa en categorías distintas validando los nombres de táctica contra la
relación oficial, y conserva las no reconocidas en un campo aparte en lugar de descartarlas.

\subsection{Resultados}

\subsubsection{Conjunto de referencia}

Sobre el conjunto atómico de robo de credenciales, el sistema produce \textbf{diez alertas y
ningún falso positivo}, procedentes de cinco reglas independientes pertenecientes a tres
categorías de origen distintas. Todas coinciden en señalar la técnica documentada en los
metadatos del conjunto de datos, lo que confirma la corrección del mapeo.

Resulta relevante que las cinco reglas detectan el mismo ataque por vías diferentes: unas por
el acceso a la memoria del proceso de credenciales, otras por la línea de comandos empleada,
otras por el fichero de volcado escrito a disco. Esta redundancia entre fuentes independientes
constituye evidencia corroborada frente a un posible falso positivo.

\subsubsection{Escenario multi-etapa}

Sobre el escenario de emulación de campaña completa, el sistema procesa 195.197 eventos
normalizados y produce \textbf{741 alertas} (0,38\,\% de los eventos), sin ningún rechazo. Las
alertas implican \textbf{cuatro equipos} distintos y cubren \textbf{33 técnicas} de ATT\&CK,
incluyendo robo de credenciales, persistencia, ejecución mediante servicios y exfiltración.

La distribución temporal es reveladora: aproximadamente el 75\,\% de las alertas se concentra
en el 20\,\% del intervalo analizado, lo que se corresponde con la ventana en que se ejecuta
el grueso de la campaña.

\subsubsection{Fatiga de alertas: evidencia empírica}

El análisis de las alertas producidas ilustra con datos propios el problema que motiva este
trabajo. Una única regla genera 292 alertas, el 39\,\% del total. Su inspección revela que
todas proceden de \textbf{una sola acción del adversario}: el despliegue de un conjunto de
herramientas administrativas mediante un intérprete de comandos, que generó una alerta por
cada uno de los ficheros extraídos, en un intervalo de ocho minutos y sobre un único equipo.

La detección es \textbf{correcta}: el comportamiento observado es efectivamente sospechoso. Lo
inadecuado es la granularidad. Un analista no necesita 292 avisos, sino uno que informe de que
se desplegó un conjunto de herramientas en un equipo determinado.

Este caso concreto cuantifica el problema: una acción del adversario consume 292 unidades de
atención del analista. Constituye asimismo el caso de uso más claro para la correlación
descrita en la Sección siguiente.

\subsubsection{Potencial de reducción de alertas}

Como referencia previa a la implementación de la correlación multi-etapa, se estimó el efecto
de una agrupación elemental por regla, equipo y ventana temporal de diez minutos. El resultado
reduce las 741 alertas a \textbf{87 grupos}, un \textbf{88,3\,\%} de reducción, con un grupo
máximo de 292 elementos.

Esta cifra constituye el \textbf{suelo} contra el que debe medirse el Pilar~A: una heurística
trivial ya alcanza ese nivel, por lo que la correlación basada en patrones debe superarlo
agrupando además alertas de reglas distintas que constituyen una misma cadena de ataque.

\subsubsection{Cobertura ATT\&CK}

Se generó la matriz de cobertura en el formato del ATT\&CK Navigator \cite{navigator}, con dos
capas: la cobertura \textit{potencial} —técnicas para las que existe al menos una regla
cargada— y la \textit{observada} —técnicas efectivamente detectadas sobre los conjuntos de
datos procesados.

\begin{table}[H]
\centering
\begin{tabular}{|l|r|}
\hline
\textbf{Indicador} & \textbf{Valor} \\
\hline
Reglas cargadas & 2.301 \\
Reglas sin etiqueta ATT\&CK & 238 (10,3\,\%) \\
Técnicas con cobertura potencial & 300 \\
\quad de las cuales subtécnicas & 177 \\
Técnicas observadas en los datos & 34 \\
Tasa de activación & 11,3\,\% \\
\hline
\end{tabular}
\caption{Cobertura ATT\&CK del sistema.}
\label{tab:cobertura-attack}
\end{table}

La baja tasa de activación no indica una limitación del sistema, sino de los conjuntos de
datos empleados: existen 266 técnicas con detección disponible que ninguno de los escenarios
procesados llega a ejercitar. Esta observación motiva la construcción de un entorno de
emulación propio que permita activar dicha cobertura latente de forma dirigida.

\subsubsection{Rendimiento}

Sobre el escenario completo, el evaluador procesa 78 eventos por segundo frente a 2.093 reglas
simultáneas, lo que equivale a aproximadamente 163.000 evaluaciones de regla por segundo. El
consumidor está implementado sobre un grupo de Kafka sin asignación fija de particiones, de
modo que el rendimiento es escalable mediante la ejecución de instancias adicionales.

\section{Detección de anomalías}
\label{sec:anomalias}

Las reglas descritas en la sección anterior detectan comportamiento \textbf{documentado}. Para
complementarlas se incorpora un detector de anomalías, cuyo cometido es señalar actividad
inusual aunque no exista regla que la describa.

\subsection{Papel dentro del sistema}

Una decisión de diseño relevante es que la anomalía \textbf{no genera alerta por sí misma}.
La razón es su elevada tasa de falsos positivos: emitir un aviso por cada desviación
estadística agravaría precisamente el problema que el trabajo pretende resolver. Su salida se
emplea en cambio para elevar la confianza de una secuencia ya detectada y para marcar
actividad que merece revisión en ausencia de regla aplicable.

\subsection{Nivel de agregación y características}

El modelo opera sobre \textbf{ventanas de actividad por equipo} en lugar de sobre eventos
individuales, por su mejor correspondencia con la naturaleza multi-etapa del problema: un
ataque se manifiesta como un patrón de actividad, no como un evento aislado.

Se definen 42 características, separadas deliberadamente en dos familias:

\begin{itemize}
    \item \textbf{Absolutas} (26): recuentos de eventos por clase, tipo de actividad y señal,
    así como medidas de diversidad. Dependen del volumen de actividad de la ventana.
    \item \textbf{Relativas} (16): proporciones de cada tipo de actividad sobre el total.
    Son \textbf{independientes del volumen por construcción}.
\end{itemize}

Esta separación no es accesoria: permite determinar experimentalmente qué parte de la
capacidad de detección procede del volumen de actividad y cuál del perfil de comportamiento.

\subsection{Protocolo experimental}

El entrenamiento emplea Isolation Forest \cite{isolationforest} sobre la implementación de
scikit-learn \cite{scikitlearn}, y se comparan dos estrategias:

\begin{enumerate}
    \item \textbf{Entrenamiento con ventanas limpias}: únicamente aquellas sin indicios
    objetivos de actividad maliciosa. Corresponde al planteamiento canónico de aprender el
    comportamiento normal.
    \item \textbf{Entrenamiento con contaminación}: todas las ventanas, declarando la
    proporción esperada de anomalías.
\end{enumerate}

Se incorpora asimismo una variante \textbf{supervisada} mediante bosques aleatorios, con y sin
compensación del desbalanceo por SMOTE \cite{smote}.

El conjunto de evaluación se mantiene \textbf{estrictamente separado} del de entrenamiento: se
reserva un escenario completo, para el que se dispone de las alertas de la capa de reglas como
referencia externa. Todos los experimentos se registran en MLflow \cite{mlflow} con sus
parámetros, métricas y semillas, garantizando su reproducibilidad.

\subsection{Controles experimentales}

Junto a los modelos se incorporan dos controles cuya función es delimitar el mérito real de
los resultados:

\begin{itemize}
    \item \textbf{Control de volumen}: un modelo entrenado con una \textbf{única}
    característica, el número de eventos de la ventana. Si alcanzase un rendimiento comparable
    al del modelo completo, quedaría demostrado que la aportación del enfoque multivariante es
    marginal.
    \item \textbf{Ablación de características potencialmente circulares}: se excluyen aquellas
    que participan en la definición de «ventana limpia», para descartar que el modelo
    aproveche información derivada de su propio criterio de entrenamiento.
\end{itemize}

\subsection{Resultados}

\begin{table}[H]
\centering
\begin{tabular}{|l|c|c|c|}
\hline
\textbf{Configuración} & \textbf{Caract.} & \textbf{AUC-ROC} & \textbf{AUC-PR} \\
\hline
No supervisado, limpias, todas & 42 & \textbf{0{,}995} & 0{,}996 \\
No supervisado, limpias, sin circulares & 34 & 0{,}995 & 0{,}996 \\
\textit{Control: solo volumen} & 1 & \textit{0{,}941} & \textit{0{,}920} \\
Supervisado sin SMOTE, sin circulares & 34 & 0{,}923 & 0{,}882 \\
Supervisado con SMOTE, solo ratios & 16 & 0{,}913 & 0{,}869 \\
No supervisado, contaminado, todas & 42 & 0{,}913 & 0{,}919 \\
No supervisado, limpias, solo ratios & 16 & 0{,}867 & 0{,}806 \\
Supervisado con SMOTE, sin circulares & 34 & 0{,}864 & 0{,}829 \\
Supervisado sin SMOTE, solo ratios & 16 & 0{,}854 & 0{,}798 \\
No supervisado, contaminado, solo ratios & 16 & 0{,}533 & 0{,}646 \\
\hline
\end{tabular}
\caption{Resultados del detector de anomalías sobre el escenario de evaluación.}
\label{tab:anomalias}
\end{table}

De la Tabla~\ref{tab:anomalias} se desprenden cuatro observaciones:

\begin{enumerate}
    \item \textbf{No existe circularidad.} Excluir las características que participan en la
    definición de ventana limpia no degrada el resultado.

    \item \textbf{Volumen y perfil son fuentes de información complementarias.} El control
    basado únicamente en el volumen alcanza 0,941, y el modelo que emplea exclusivamente
    características relativas —independientes del volumen— alcanza 0,867. Ambas discriminan de
    forma independiente, y su combinación (0,995) supera a cualquiera por separado, lo que
    justifica el enfoque multivariante frente a una regla de umbral.

    \item \textbf{La pureza del conjunto de entrenamiento resulta determinante.} Con idéntico
    conjunto de características, entrenar con ventanas limpias produce 0,867 mientras que
    hacerlo con datos contaminados produce 0,533, indistinguible del azar.

    \item \textbf{El efecto de SMOTE depende de la dimensionalidad.} Mejora el resultado con
    16 características (0,854 $\rightarrow$ 0,913) y lo degrada con 34 (0,923 $\rightarrow$
    0,864), comportamiento coherente con la interpolación entre vecinos en espacios de
    dimensión elevada.
\end{enumerate}

\subsection{Limitaciones}

Los resultados deben interpretarse con dos cautelas explícitas.

La primera es el \textbf{tamaño del conjunto de evaluación}: veintiocho ventanas, de las cuales
quince positivas. Con esa magnitud, la incertidumbre asociada a las métricas es
considerable y valores próximos a la unidad deben leerse como indicio de separabilidad
sencilla antes que como evidencia de un modelo excelente.

La segunda, más relevante, afecta a la \textbf{representatividad de los datos}. Las ventanas
que contienen actividad maliciosa presentan una mediana de eventos ocho veces superior a las
benignas, con solapamiento reducido. Esta separabilidad es en buena medida un artefacto de la
naturaleza de los conjuntos de datos empleados: capturas de laboratorio en las que las
técnicas se ejecutan de forma concentrada sin actividad de fondo sostenida. En un entorno
real, la actividad legítima presenta volúmenes comparables y un adversario competente opera
precisamente con discreción.

Ambas limitaciones convergen en la misma necesidad: disponer de telemetría con actividad
benigna prolongada y de volumen comparable al de la actividad maliciosa. Esta constatación
motiva la construcción del entorno de emulación descrito a continuación.

\section{Laboratorio de emulación propio}
\label{sec:laboratorio}

Para disponer de telemetría con actividad benigna sostenida y etiquetado exacto se construyó
un laboratorio aislado: una máquina virtual Windows~11 en VirtualBox, con Sysmon y la
configuración de referencia de la comunidad, y Atomic Red Team para ejecutar técnicas
ATT\&CK concretas. El mecanismo de captura registra la ventana temporal de cada ejecución en
un manifiesto, exporta los eventos al terminar y etiqueta cada uno de forma diferida por
pertenencia a una ventana, sin búsqueda de texto en vivo.

El conjunto final contiene 14.895 eventos, 14.765 benignos procedentes de tres horas de
actividad simulada y 130 maliciosos de cuatro técnicas (\texttt{T1547.001},
\texttt{T1003.001}, \texttt{T1055} y \texttt{T1070}). Las técnicas que Windows Defender o la
protección de LSA bloquearon se documentan como «intentadas y prevenidas por un control», sin
forzarlas, y se descartó expresamente el uso de herramientas de malware real.

La revisión del primer conjunto destapó tres \textbf{fuentes de contaminación} que habrían
invalidado el etiquetado sin producir error alguno:

\begin{itemize}
    \item El registro de módulos de PowerShell suponía el 59\,\% del conjunto y era casi
    íntegramente ruido interno del propio marco de pruebas, repartido por igual entre ambas
    clases. Se desactivó y se conservó solo el registro de bloques de guion, que captura el
    contenido real.
    \item Las llamadas de contexto internas del ejecutor de pruebas (\texttt{whoami.exe},
    \texttt{hostname.exe}) suponían el 62\,\% de las muestras maliciosas a nivel de proceso.
    Tras filtrarlas, las muestras maliciosas pasaron de 1.052, con una sola ventana que
    concentraba el 78\,\%, a 147 bien repartidas entre técnicas.
    \item Un servicio del sistema escribía telemetría de compatibilidad de aplicaciones
    durante las ventanas de ataque, lo que añadía 16 filas mal etiquetadas.
\end{itemize}

La integración de este conjunto en la evaluación requiere una ruta de conversión al esquema
OCSF que todavía no existe, y queda planificada para la fase de evaluación.

\section{Pilar A: detección de secuencias multi-etapa}
\label{sec:pilarA}

La detección basada en reglas descrita en la Sección~\ref{sec:deteccion-reglas} produce
alertas aisladas: cada una señala una técnica concreta en un instante concreto. El Pilar~A
transforma esas alertas en \textbf{campañas}, es decir, en conjuntos de evidencias que el
sistema reconoce como etapas sucesivas de un mismo ataque. Esta sección describe el corpus
sobre el que se ha construido, el método de validación, la estructura de correlación en dos
niveles que impone la herramienta, los tres patrones implementados, el esquema común en que
se emiten las campañas y su validación en modo continuo.

\subsection{Corpus de evaluación y verdad conocida}

Todo el trabajo de esta fase se apoya en un \textbf{corpus fijo y reproducible}, cargado
mediante un único guion que recrea los \textit{topics}, reproduce los conjuntos de datos y
los procesa con grupos de consumo nuevos (Anexo~\ref{anx:reproducibilidad}). El corpus reúne
nueve conjuntos de datos: el primer día de la emulación de APT29 publicada por MITRE, una
captura atómica de volcado de LSASS mediante \texttt{comsvcs.dll} y siete campañas de robo
de credenciales que emplean siete herramientas distintas.

\begin{table}[H]
\centering
\begin{tabular}{|l|r|}
\hline
\textbf{Magnitud} & \textbf{Valor} \\
\hline
Eventos consumidos & 550.494 \\
Eventos normalizados a OCSF & 530.401 (96,4\,\%) \\
Eventos derivados a la cola de rechazos & 20.093 \\
Reglas Sigma evaluadas & 2.093 \\
Alertas distintas producidas & 1.184 \\
\quad de ellas, en el primer día de APT29 & 741 \\
\hline
\end{tabular}
\caption{Corpus de evaluación del Pilar A.}
\label{tab:corpus-f3}
\end{table}

La totalidad de los rechazos corresponde a tipos de evento sin mapeador, y en torno al 60\,\%
son marcadores de inicio y fin de bloques de PowerShell sin contenido. Ningún evento de
acceso a procesos ni de creación de ficheros figura entre ellos, de modo que la evidencia
central del robo de credenciales se normaliza sin pérdida.

La \textbf{verdad conocida} del primer día de APT29 se reconstruyó a partir del plan de
emulación publicado: 28 pasos agrupados en 10 etapas, con su equipo y usuario. Los
identificadores de técnica del plan corresponden a versiones antiguas de ATT\&CK y varios
están obsoletos (por ejemplo, \texttt{T1086} es hoy \texttt{T1059.001} y \texttt{T1088} es
\texttt{T1548.002}). Sin traducirlos, el cruce entre el plan y las alertas no cuadra y
aparenta fallos de detección que no lo son.

\subsection{Método: versión de referencia y versión en flujo}

Cada patrón se ha implementado \textbf{dos veces}. La primera es una vista SQL sobre el
almacén analítico que procesa el corpus completo por lotes y actúa como
\textbf{oráculo}: fija el resultado esperado. La segunda es la implementación en Flink con
\texttt{MATCH\_RECOGNIZE}, que debe reproducir ese resultado.

El valor de este método es separar las fuentes de error. Sin oráculo, depurar el patrón
obliga a dudar a la vez de su lógica y de la semántica temporal del motor; con él, la
pregunta se reduce a por qué no aparece lo que se sabe que existe. El método funciona en
ambos sentidos: en un caso fue la versión en flujo la que reveló un defecto del oráculo, que
descartaba la captura atómica por exigir contexto alrededor de las evidencias, requisito que
penaliza precisamente las capturas más limpias.

Para validar en Flink sin depender del avance de las marcas de agua, la fuente del bus se
configura como \textbf{acotada} (se lee hasta el último desplazamiento disponible) y el
motor se ejecuta en modo por lotes. Así cada consulta produce un resultado final, comparable
fila a fila con el oráculo. La validación en modo continuo se aborda por separado en la
Sección~\ref{sec:pilarA-streaming}.

\subsection{Correlación en dos niveles}
\label{sec:dos-niveles}

La intuición inicial era expresar cada campaña con un único patrón: «todas las evidencias
consecutivas de un mismo equipo hasta que se abra un hueco de más de diez minutos». Esa
expresión resulta \textbf{imposible} con \texttt{MATCH\_RECOGNIZE} en Flink. El motor no
admite cuantificadores voraces como último elemento de un patrón, y el cuantificador
reluctante, que sí admite, produce la coincidencia \textit{mínima}: con
\texttt{PATTERN (E+? C)} se obtienen 29 coincidencias de dos filas donde el oráculo
reconoce 10 campañas.

La solución es una \textbf{correlación jerárquica en dos niveles}, que aquí deja de ser una
opción de diseño para convertirse en un requisito impuesto por la herramienta:

\begin{enumerate}
    \item \textbf{Nivel 1, transiciones.} \texttt{MATCH\_RECOGNIZE} reconoce pares de
    evidencias consecutivas que cumplen la condición del patrón dentro de una ventana
    temporal. Es lo que un motor CEP hace bien.
    \item \textbf{Nivel 2, campañas.} Una agregación agrupa las transiciones separadas por
    un hueco menor que la ventana. Es lo que un motor CEP hace mal y una agregación resuelve
    de forma trivial. En lotes se implementa con funciones de ventana analíticas; en flujo,
    con una ventana de sesión sobre el tiempo de cada transición.
\end{enumerate}

\begin{figure}[H]
\centering
\begin{tikzpicture}[
  ev/.style={circle,draw=black!70,fill=black!10,minimum size=4mm,inner sep=0pt},
  evc/.style={ev,fill=black!55},
  tr/.style={-{Stealth[length=2mm]},thick},
  font=\footnotesize]
  \foreach \x/\t in {0/ev,1.1/evc,2.0/ev,3.2/evc,7.4/ev,8.4/evc,9.3/ev} {
    \node[\t] at (\x,0) {};
  }
  \draw[tr] (0.15,0.25) to[bend left=40] (0.95,0.25);
  \draw[tr] (1.25,0.25) to[bend left=40] (1.85,0.25);
  \draw[tr] (2.15,0.25) to[bend left=40] (3.05,0.25);
  \draw[tr] (7.55,0.25) to[bend left=40] (8.25,0.25);
  \draw[tr] (8.55,0.25) to[bend left=40] (9.15,0.25);
  \node[anchor=west] at (9.8,0.45) {Nivel 1: transiciones};
  \draw[thick] (-0.2,-0.45) -- (3.4,-0.45);
  \draw[thick] (-0.2,-0.35) -- (-0.2,-0.55); \draw[thick] (3.4,-0.35) -- (3.4,-0.55);
  \draw[thick] (7.2,-0.45) -- (9.5,-0.45);
  \draw[thick] (7.2,-0.35) -- (7.2,-0.55); \draw[thick] (9.5,-0.35) -- (9.5,-0.55);
  \node at (1.6,-0.8) {campaña 1};
  \node at (8.35,-0.8) {campaña 2};
  \draw[<->,dashed] (3.4,-0.15) -- node[above] {hueco $>$ ventana} (7.2,-0.15);
  \node[anchor=west] at (9.8,-0.45) {Nivel 2: agregación por hueco};
  \draw[->] (-0.4,-1.25) -- (12.8,-1.25) node[right] {$t$};
\end{tikzpicture}
\caption{Correlación en dos niveles. Los círculos oscuros son evidencias de la etapa ancla
del patrón; las flechas, las transiciones reconocidas por \texttt{MATCH\_RECOGNIZE}.}
\label{fig:dos-niveles}
\end{figure}

La implementación encontró además otras restricciones del motor, que condicionan la forma
de escribir los patrones: la cláusula \texttt{DEFINE} no admite subconsultas, por lo que la
lógica sobre listas de técnicas se precalcula en la proyección; \texttt{COUNT(DISTINCT)} no
está permitido en \texttt{MEASURES}; y \texttt{WITHIN} acota la duración total de la
coincidencia y no el hueco entre elementos, por lo que la ventana se expresa como condición
en \texttt{DEFINE}.

Una restricción adicional determina la arquitectura. \texttt{MATCH\_RECOGNIZE} exige ordenar
por un atributo de tiempo de evento, y Flink lo pierde al atravesar uniones y agregaciones.
Por ello, cada patrón se divide en dos trabajos desacoplados por el bus: uno de
\textbf{preparación} publica en un \textit{topic} intermedio la evidencia del patrón, sin
repeticiones y en orden de tiempo de evento, y otro de \textbf{reconocimiento} la relee con
su marca de agua. Es también la arquitectura adecuada para producción, pues separa el
enriquecimiento del reconocimiento. El Anexo~\ref{anx:patrones} recoge el código del
patrón P1 como ejemplo.

\subsection{P1: robo de credenciales}

P1 reconoce la secuencia herramienta $\rightarrow$ acceso a la memoria de LSASS
$\rightarrow$ volcado $\rightarrow$ limpieza. Su diseño se decidió con datos en tres puntos.

\textbf{El ancla es la familia \texttt{T1003.*} y no la subtécnica \texttt{T1003.001}.}
Medido sobre las siete variantes, la subtécnica exacta cubre seis: falla precisamente
Mimikatz, la variante más emblemática, porque Meterpreter lo ejecuta dentro de su propio
proceso y ningún proceso externo abre la memoria de LSASS. La familia cubre las siete.

\textbf{Las etapas se reconocen por el tipo de evento de origen}, restringido a alertas de
la familia \texttt{T1003}: un acceso a proceso es acceso a LSASS, una creación de proceso es
ejecución de la herramienta y una creación de fichero es volcado a disco. Una primera versión
reconocía las etapas por los campos coincidentes de la regla y producía evidencia falsa: las
292 alertas de un \textit{dropper} de PowerShell de APT29 satisfacían la etapa «volcado a
fichero» sin relación alguna con credenciales. Con la restricción, la evidencia de P1 en el
equipo SCRANTON pasa de 308 alertas a 3.

\textbf{La transición es bidireccional}: basta con que una de sus dos evidencias sea de
credenciales, de modo que valen tanto «preparación $\rightarrow$ credenciales» como
«credenciales $\rightarrow$ limpieza». Con la versión unidireccional se perdía el equipo
NASHUA, donde la limpieza llega después del único acceso a credenciales. La ventana es de
10 minutos, el doble del núcleo observado en APT29 (4\,min 47\,s entre la descarga de la
herramienta y el acceso a LSASS).

La confianza se \textbf{descompone en factores}: 0,40 por el ancla, más 0,25 si hay acceso a
LSASS, 0,20 si hay volcado a disco, 0,10 si hay ejecución de la herramienta y 0,15 si hay
limpieza, con un máximo de 1. El resultado son \textbf{11 campañas}: una por cada una de
las \textbf{siete variantes} de robo de credenciales, dos en SCRANTON y una en NASHUA dentro
de APT29, y la captura atómica. La confianza oscila entre 0,50 (Mimikatz, solo ancla y
contexto) y 0,95 (las variantes con acceso a LSASS y fichero de volcado).

\subsection{P2: movimiento lateral}

P2 reconoce el salto de un equipo a otro: una conexión de red hacia el destino seguida de
ejecución remota en él. Es el primer patrón que \textbf{cruza equipos}, y requiere
construir antes las \textbf{aristas} de red a partir de los eventos normalizados.

La construcción de las aristas destapó varios problemas que no producían error alguno.
Bajo la misma clase OCSF de actividad de red conviven las conexiones lógicas de Sysmon y
cada flujo permitido por el cortafuegos de Windows; filtrar por clase daba 3.480 conexiones
WinRM donde el fichero de origen tiene 12. El mapa de direcciones IP a equipos, si cuenta
también los registros entrantes, se corrompe en silencio: un equipo llegó a ganar su propia
dirección por un solo voto. Se resolvió contando solo conexiones salientes, en las que la
dirección de origen es por definición la propia. Y la dirección de bucle local escrita sin
abreviar fabricaba aristas falsas hasta que se normalizó su grafía.

Con el mapa corregido y el controlador de dominio identificado por su comportamiento y
excluido, quedan 39 aristas: 22 hacia el controlador y \textbf{17 de SCRANTON a NASHUA}
(puertos 135, 445 y 5985), todas confirmadas desde los dos extremos. Sobre ellas, el nivel~1
reconoce en cada destino una arista y una evidencia de ejecución (técnicas \texttt{T1021},
\texttt{T1569.002} y \texttt{T1003}) separadas como máximo 300 segundos, en cualquier orden,
y el nivel~2 agrupa las transiciones de un mismo par de equipos.

El resultado es \textbf{una campaña}, idéntica campo a campo a la del oráculo: de las
03:09:23.801 a las 03:17:59.493, con 8 evidencias en el destino, cuatro etapas (conexión,
ejecución remota, credenciales y servicio) y confianza 1,0. El orden no se exige pero se
\textbf{mide}: siete de las ocho evidencias llegan después de su conexión, con un retraso
máximo de 1.004\,ms. La octava tiene una hora de resolución de segundo y no permite ordenar
dentro de ese segundo.

\subsection{P3: escalada de privilegios con descubrimiento o persistencia}

El inventario previo mostró que \textbf{las reglas no ven el descubrimiento del atacante}
en APT29: solo hay dos alertas de la táctica en todo el día, y ambas proceden del agente de
la nube del equipo. El descubrimiento real está en los bloques de PowerShell registrados
(\texttt{Invoke-Discovery}, \texttt{Get-Process} remoto). Por ello P3 es el primer patrón
\textbf{híbrido}: su etapa de descubrimiento combina alertas Sigma y eventos de PowerShell
con instrucciones de reconocimiento.

El ancla es la \textbf{escalada de privilegios} (técnicas \texttt{T1548}, \texttt{T1134},
\texttt{T1055} y \texttt{T1068}), acompañada de descubrimiento o de persistencia
(\texttt{T1547}, \texttt{T1543}, \texttt{T1053}, \texttt{T1546}, \texttt{T1037} y
\texttt{T1136}) en el mismo equipo. Los eventos de asignación de privilegios especiales se
descartaron como etapa porque aparecen en todos los equipos de forma rutinaria. Se aplica
una única exclusión, medida: la modificación de claves de arranque que escribe el
explorador de Windows al iniciar sesión, que eliminaba 28 evidencias espurias. Una segunda
exclusión prevista para los agentes de la nube se retiró al comprobar que no eliminaba
nada: una exclusión sin efecto es un ajuste especulativo.

El resultado son \textbf{3 campañas}. En APT29, la de SCRANTON (confianza 0,90) reproduce
el orden exacto del plan del atacante: escalada $\rightarrow$ descubrimiento $\rightarrow$
persistencia. Las otras dos son las capturas de robo de credenciales que incluyen escalada
real (confianza 0,75). Lo que el patrón \textbf{deja fuera} es tan relevante como lo que
detecta: el controlador de dominio, con 25 evidencias de descubrimiento producidas por un
guion de monitorización legítimo, y cinco capturas con entre 26 y 64 evidencias de
reconocimiento no forman campaña porque no hay escalada. El software legítimo hace el mismo
reconocimiento que un atacante; es la secuencia lo que los distingue.

\begin{table}[H]
\centering
\small
\begin{tabular}{|l|p{4.3cm}|p{5.2cm}|}
\hline
\textbf{Patrón} & \textbf{Qué reconoce} & \textbf{Resultado} \\
\hline
P1 Robo de credenciales & Herramienta, acceso a LSASS, volcado y limpieza & 11 campañas; las 7 variantes de robo de credenciales detectadas \\
\hline
P2 Movimiento lateral & Conexión a otro equipo y ejecución remota en el destino & 1 campaña SCRANTON $\rightarrow$ NASHUA, 17 aristas bilaterales \\
\hline
P3 Escalada & Escalada con descubrimiento o persistencia & 3 campañas; orden del plan de APT29 reproducido \\
\hline
\end{tabular}
\caption{Patrones del Pilar A y resultados sobre el corpus. En los tres, la versión en
Flink reproduce exactamente su oráculo.}
\label{tab:patrones}
\end{table}

\subsection{Esquema común de campaña}

Las campañas de los tres patrones se emiten con un \textbf{esquema común}, que es lo que
consumirán el grafo de ataque, la priorización y la interfaz. Por coherencia con el resto
del sistema, una campaña es un hallazgo estándar \textbf{OCSF 1.5.0 «Detection Finding»}
(clase 2004). Todo lo que el estándar contempla va en sus campos: las alertas relacionadas,
las técnicas, la confianza y el intervalo temporal. Lo que no contempla va en una extensión
propia documentada: las \textbf{etapas ordenadas}, los \textbf{factores de confianza}, cuya
suma es la confianza publicada, los equipos con su papel y las métricas de cada patrón. Las
técnicas se publican sin táctica asociada, porque las alertas traen ambas listas sin
emparejar y emparejarlas sería inventar.

Dos reglas garantizan que el resultado sea reproducible. La identidad de cada campaña es un
resumen criptográfico de su conjunto de datos, patrón, clave y hora de inicio, de modo que
reprocesar el corpus produce los mismos identificadores y la tabla de destino sobrescribe en
lugar de duplicar. Las etapas se ordenan por el segundo de su primera evidencia y, dentro
del mismo segundo, por el orden lógico del patrón, porque parte de las horas de origen solo
tienen resolución de segundo. El proceso que vuelca las campañas al almacén
\textbf{valida cada documento} y se detiene ante el primero que no cumpla el esquema. El
almacén contiene hoy 15 campañas: 11 de P1, 1 de P2 y 3 de P3. El
Anexo~\ref{anx:esquema} detalla el esquema con un ejemplo completo.

\subsection{Validación en modo continuo}
\label{sec:pilarA-streaming}

La validación anterior se realizó con fuentes acotadas. En modo continuo, la marca de agua
descarta sin aviso los eventos que llegan con más retraso del admitido, por lo que antes de
lanzar los patrones se \textbf{midió el desorden} de llegada:

\begin{itemize}
    \item Dentro de un mismo conjunto de datos, las alertas llegan con hasta 293\,s de
    retraso respecto a otras posteriores, por efecto de los detectores en paralelo. Con una
    marca de agua de 30\,s se perderían 220 alertas (18\,\%); con 5 minutos, ninguna. La
    cifra está inflada porque el corpus se reproduce decenas de veces más rápido que el
    tiempo real.
    \item Los \textit{topics} intermedios, escritos inicialmente sin ordenar, habrían
    llegado «tarde» en un 74--78\,\%. Se decidió publicarlos en orden de tiempo de evento.
    \item Al reproducir a la vez capturas de distintos años, el desorden entre conjuntos de
    datos alcanza años. Es un artefacto de la reproducción que no existe en producción.
\end{itemize}

Con los \textit{topics} intermedios ordenados, el segundo nivel se expresa con una ventana de
sesión y los tres patrones dan en flujo \textbf{exactamente el mismo resultado} que por lotes:
las 11 campañas de P1 con sus 52 transiciones, la campaña de P2 y las 3 de P3. La ventana de
sesión no emite cada campaña una sola vez al cerrarse, sino que la \textbf{actualiza} con cada
transición nueva. Para detectar es una ventaja, porque la campaña existe desde su primera
transición, pero obliga a emitirla hacia un destino que admita actualizaciones.

\subsection{Incidencia: la hora de los eventos era la de recogida}

El primer día de APT29 no incluye el campo de hora de creación del evento, y la
normalización recurría a la marca del colector, que corresponde al instante de
\textbf{recepción}. El desfase medido sobre los eventos que sí traen la hora real tiene una
mediana de 6,2\,s, un percentil 95 de 74,9\,s y un máximo de 101,2\,s. Lo grave es su
dispersión, de 26,7\,s: los eventos se reordenaban entre sí y no había desplazamiento
constante que corregir.

La corrección establece una cadena de prioridad entre las fuentes de tiempo disponibles, con
la conversión de zona horaria derivada por conjunto de datos, y marca cada evento con la
fuente empleada. Ninguna fuente es fiable en todos los conjuntos: la hora UTC de Sysmon, la
más precisa en APT29, va desviada 19 horas en las capturas de robo de credenciales. Con la
hora corregida desapareció una inversión imposible documentada hasta entonces, en la que el
proceso remoto aparecía tres segundos antes que la conexión que lo provoca: la conexión va
ahora 116\,ms antes. Queda una limitación del dato de origen: el 26,6\,\% de los eventos de
APT29 solo tiene resolución de segundo.

\subsection{Incidencia: la identidad de los eventos y de las alertas}

Durante esta fase se comprobó que el 15\,\% de los eventos normalizados de APT29 era idéntico
byte a byte a otro, porque el evento OCSF no conservaba el número de registro de origen. Dos
ocurrencias distintas con los mismos campos se volvían indistinguibles, y las alertas
duplicadas por un reprocesado no se podían separar de las legítimas. Se resolvió dando
\textbf{identidad determinista} a ambos: el evento conserva su número de registro (o, si la
fuente no lo aporta, el ordinal de la línea en el fichero), la alerta se identifica por el
resumen del evento y la regla, y las tablas del almacén deduplican por esa identidad. La
comprobación final fue recuperar exactamente las 741 alertas documentadas para APT29 y las
cifras de eventos de cada conjunto de datos contadas contra su fichero de origen.

\subsection{Estado de la fase y trabajo pendiente}

Se cumplen los criterios de cierre relativos a la detección: tres patrones que disparan sobre
datos reales, las siete variantes de robo de credenciales detectadas y campañas multi-etapa
reconocidas en APT29, con resultados idénticos por lotes y en flujo. Quedan pendientes:

\begin{itemize}
    \item Medir los \textbf{falsos positivos} sobre la porción benigna del corpus y sobre
    capturas atómicas de una sola táctica.
    \item Medir la \textbf{ratio de reducción de alertas} de los tres patrones. La de P1
    (0,991) no es comparable con la línea base de agrupación de la fase anterior, porque solo
    cubre lo que P1 absorbe alrededor de sus anclas.
    \item Decidir la \textbf{extensión de una campaña}: la versión en Flink la acota a su
    evidencia y el oráculo absorbe $\pm$10 minutos. En P3 la diferencia es grande, porque el
    descubrimiento llega en ráfagas alrededor del ancla, y condiciona la ratio de reducción.
    \item Resolver el \textbf{punto ciego del etiquetado}: el 27,1\,\% de las alertas no lleva
    técnica ATT\&CK y es invisible para la correlación por construcción.
    \item Incorporar el segundo día de APT29 y emitir las campañas en flujo hacia un destino
    que admita actualizaciones.
\end{itemize}


% (oculta hasta F4-F6) \section{Plano de inteligencia de amenazas}
% TODO (F4): despliegue de la plataforma CTI e importación de ATT&CK; sincronización de
% fuentes de indicadores; extracción automática de técnicas a partir de informes en lenguaje
% natural y carga del resultado en el grafo.

% (oculta hasta F4-F6) \section{Pilar B: grafo de ataque y razonamiento}
% TODO (F5): modelo de datos del grafo; poblado del inventario y las vulnerabilidades;
% analítica de grafos para criticidad y caminos; reglas de interacción en Datalog; exportación
% de hechos y reimportación de las rutas derivadas; trazabilidad de la derivación.

% (oculta hasta F4-F6) \section{Pilar C: priorización explicable y servicios}
% TODO (F6): modelo de puntuación y justificación de los pesos; análisis de sensibilidad;
% generación de la narrativa; interfaz de programación de la aplicación.


\chapter{Resultados}
\label{cap:resultados}

Este capítulo recogerá la evaluación cuantitativa del sistema completo, que se realizará en
la fase de evaluación una vez construidos los tres pilares. Los resultados parciales obtenidos
hasta ahora en cada componente se presentan junto a su desarrollo en el
Capítulo~\ref{cap:desarrollo}; en particular, la reducción de alertas y los falsos positivos
del Pilar~A se miden en la Sección~\ref{sec:pilarA-evaluacion}. El protocolo se fija \textbf{antes} de ejecutar los
experimentos, para evitar el sesgo de elegir las métricas una vez conocidos los resultados.

\section{Protocolo de evaluación}

\subsection{Conjunto de evaluación y verdad conocida}

La evaluación empleará conjuntos de datos públicos cuya actividad maliciosa está documentada:
las siete campañas de robo de credenciales, los dos días de la emulación de APT29, capturas de
explotación de vulnerabilidades y de falsificación de credenciales federadas, y las capturas
atómicas de host de OTRF, de las que el Pilar~A ya emplea 101. La verdad conocida de cada conjunto se consolidará en
metadatos estructurados (técnicas, equipos, ventanas temporales y, cuando exista, el plan de
emulación), de modo que todas las métricas se calculen contra la misma referencia.

\subsection{Familias de métricas}

\begin{itemize}
    \item \textbf{Detección}: precisión, exhaustividad y F1 por técnica y en conjunto; tasa de
    falsos positivos sobre tráfico benigno etiquetado, que las capturas atómicas no aportan
    porque son todas maliciosas; cobertura de ATT\&CK.
    \item \textbf{Correlación multi-etapa}: campañas correctamente reconstruidas, completitud
    de sus etapas y ratio de reducción de alertas.
    \item \textbf{Priorización}: presencia del incidente real entre los $K$ primeros,
    rango recíproco medio y correlación entre la puntuación y la gravedad real.
    \item \textbf{Razonamiento}: coincidencia entre las rutas derivadas y la ruta del plan de
    emulación.
    \item \textbf{Rendimiento}: rendimiento por etapa, latencia de detección (percentiles 50,
    95 y 99), consumo de recursos y comportamiento con la carga multiplicada por 10 y por 100.
    \item \textbf{Componentes de aprendizaje}: AUC-ROC y AUC-PR del detector de anomalías;
    precisión y exhaustividad por técnica del extractor, comparando el modelo base con el
    afinado; precisión de la predicción de enlaces en el grafo; mejora del orden de la
    priorización frente a pesos fijos; y coste computacional de cada uno.
\end{itemize}

\subsection{Líneas base}

Cada aportación se medirá frente a una línea base que la aísla:

\begin{table}[H]
\centering
\begin{tabular}{|c|p{6cm}|p{5cm}|}
\hline
\textbf{Línea base} & \textbf{Configuración} & \textbf{Qué demuestra} \\
\hline
B1 & Solo reglas Sigma, sin correlación & El valor del CEP \\
B2 & Solo reglas, sin detección de anomalías & El valor de la detección de anomalías \\
B3 & Priorización por la gravedad de la regla & El valor de la puntuación contextual \\
B4 & Razonamiento sin el modelo de grafo & El valor del aprendizaje sobre el grafo \\
B5 & Las mismas reglas en un SIEM abierto & Una comparación externa (opcional) \\
\hline
\end{tabular}
\caption{Líneas base de la evaluación.}
\label{tab:lineas-base}
\end{table}

\subsection{Reproducibilidad}

Toda la evaluación se ejecutará con un único comando que produzca las métricas y las figuras.
Cada experimento se repetirá al menos tres veces y se informará de la media y la desviación.
El análisis incluirá de forma explícita los casos en que el sistema falla y su explicación.

% Secciones ocultas hasta la fase de evaluación (F7):
% \section{Métricas de detección}
% \section{Métricas de correlación multi-etapa}
% \section{Métricas de priorización}
% \section{Métricas de razonamiento}
% \section{Rendimiento}
% \section{Discusión}


\chapter{Conclusiones y trabajo futuro}
\label{cap:conclusiones}

Las conclusiones y las limitaciones se redactarán al cierre del desarrollo, cuando pueda
valorarse el grado de cumplimiento de cada objetivo específico. Las líneas de trabajo futuro
identificadas hasta ahora se recogen a continuación.

% Secciones ocultas hasta el cierre del desarrollo:
% \section{Conclusiones}
% TODO: grado de cumplimiento de cada objetivo específico planteado en la introducción;
% contribuciones efectivamente aportadas; lecciones aprendidas del desarrollo.
% \section{Limitaciones}
% TODO: alcance de la validación, dependencia de los conjuntos de datos empleados,
% aspectos no cubiertos por el prototipo.

\section{Líneas de trabajo futuro}

\begin{itemize}
    \item \textbf{Ampliación de la detección.} Incorporar telemetría de red ya disponible en
    los conjuntos de datos, fuentes de Linux, nuevas familias de ataque y simulaciones propias
    de campañas completas en el laboratorio aislado. Además de ampliar la cobertura, mostraría
    que la arquitectura generaliza: una fuente nueva solo requiere un mapeador al esquema OCSF
    y sus reglas. Esta línea está planificada como ampliación condicionada a que el núcleo del
    trabajo esté cerrado a finales de febrero.
    \item \textbf{Señuelos.} Desplegar credenciales y recursos señuelo en la red real, con una
    tasa de falsos positivos casi nula. Una credencial señuelo empleada tras un volcado de LSASS
    reforzaría el patrón de robo de credenciales. Los señuelos expuestos a Internet muestran
    atacantes oportunistas y no cómo se atacaría una organización concreta, que es lo que
    aporta el análisis de rutas, y los de alta interacción plantean riesgos técnicos y legales.
    \item \textbf{Cadena de suministro de software}: extender el razonamiento sobre rutas de
    ataque a las dependencias del software desplegado.
    \item \textbf{Generación asistida del grafo}: obtener hechos del grafo de ataque a partir
    de informes en lenguaje natural.
    \item \textbf{Validación en producción}: evaluar el sistema sobre telemetría real de una
    organización, con su volumen y su actividad legítima de fondo.
\end{itemize}


%%%%% ESTILO DEL DOCUMENTO %%%%%%%


%%%%%%%%%%%%%%%%%%%%%%%%%%%%%%%%%%%%%%

%% Bibliografía. Usa el ref.bib 
\clearpage

%\nocite{*}

\Urlmuskip=0mu plus 1mu\relax

\phantomsection\addcontentsline{toc}{chapter}{Bibliografía}
\bibliography{ref.bib}
\bibliographystyle{ieeetr}

%% Anexos (no cuentan para la extensión del cuerpo)
\clearpage
\appendix
\phantomsection\addcontentsline{toc}{chapter}{Anexos}
% En los anexos, el título lleva delante «Anexo A.», «Anexo B.»...
\makeatletter
\renewcommand*\@makechapterhead[1]{
    {\Huge\bfseries \appendixname~\thechapter. #1\par \nobreak\vskip 24\p@}
}
\makeatother
\chapter{Aspectos éticos, sociales, económicos y ambientales}
\label{cap:aspectos}

\section{Introducción}

Un sistema de detección de intrusiones opera necesariamente sobre registros de la actividad
de personas y sistemas, y produce salidas que orientan decisiones con consecuencias
operativas. Ambas circunstancias exigen una reflexión explícita sobre las implicaciones
éticas y legales del trabajo, que se aborda en este anexo.

\section{Protección de datos personales}

La telemetría de seguridad contiene de forma inherente información susceptible de constituir
dato personal en el sentido del Reglamento General de Protección de Datos: nombres de
usuario, identificadores de cuenta, nombres de equipo, direcciones IP y, de forma indirecta,
el patrón de actividad de personas identificables.

En el ámbito de este trabajo, la exposición se limita de tres maneras. En primer lugar, los
datos empleados proceden exclusivamente de \textbf{conjuntos públicos capturados en entornos
de laboratorio} construidos al efecto, en los que las cuentas y equipos son ficticios y no
corresponden a personas reales. En segundo lugar, no se ha procesado telemetría de ningún
sistema en producción, propio o de terceros. Y en tercer lugar, los datos permanecen
localmente en el entorno de desarrollo, sin transferencia a servicios externos.

De cara a un eventual despliegue en un entorno real, el sistema plantearía las obligaciones
habituales de esta categoría de tratamiento: existencia de base jurídica —típicamente el
interés legítimo en la seguridad de las redes y la información—, información previa a las
personas afectadas, limitación de la finalidad al ámbito estrictamente de seguridad,
minimización de los datos recogidos, establecimiento de plazos de conservación proporcionados
y control de acceso a la información. El diseño del sistema facilita varias de estas
obligaciones: el particionado temporal del almacén permite implementar políticas de
conservación por rangos, y la separación entre campos promocionados y evento completo
posibilitaría aplicar seudonimización selectiva.

Merece mención específica el principio de \textbf{minimización} frente a la tentación
contraria propia del dominio: la práctica habitual en detección es recoger la máxima
telemetría posible por si resultara necesaria durante una investigación. Existe aquí una
tensión real entre eficacia de detección y proporcionalidad del tratamiento, que en un
despliegue real debería resolverse de forma explícita y documentada, no por omisión.

\section{Uso responsable del conocimiento ofensivo}

El trabajo emplea conocimiento sobre técnicas adversarias reales: cómo se extraen
credenciales de la memoria, cómo se progresa lateralmente por una red o cómo se eliminan las
trazas de una intrusión. Este conocimiento es intrínsecamente de doble uso.

Tres consideraciones acotan el riesgo. La primera es que todo el material empleado procede de
\textbf{fuentes públicas y ampliamente documentadas}: MITRE ATT\&CK, los conjuntos de datos
abiertos utilizados y el repositorio comunitario de reglas de detección son recursos
publicados precisamente con finalidad defensiva. El trabajo no genera conocimiento ofensivo
nuevo ni descubre técnicas no documentadas. La segunda es que la orientación del sistema es
exclusivamente \textbf{detectiva}: el objetivo es reconocer estas técnicas cuando ocurren, no
facilitar su ejecución. La tercera es que la emulación de comportamiento adversario, prevista
en fases posteriores, se realizará únicamente sobre entornos de laboratorio aislados y
controlados.

\section{Sesgo, fiabilidad y supervisión humana}

Un sistema que prioriza incidentes influye directamente en qué recibe atención y qué no. Un
sesgo en la priorización no produce un error visible: produce la \textbf{ausencia silenciosa}
de una alerta que sí debería haberse atendido. Esta consideración ha determinado dos
decisiones de diseño con dimensión ética además de técnica.

La primera es la \textbf{explicabilidad} como requisito, no como añadido. Una puntuación
opaca traslada al analista la disyuntiva de confiar ciegamente o desconfiar por completo;
una puntuación descompuesta en factores interpretables le permite evaluar críticamente el
juicio del sistema y detectar cuándo su razonamiento no aplica al caso concreto.

La segunda es el carácter de \textbf{soporte a la decisión, y no de decisión automatizada}.
El sistema no ejecuta acciones de respuesta: no aísla equipos, no bloquea cuentas ni
interrumpe conexiones. Ordena y justifica, mientras que la decisión permanece en el analista.
Esta delimitación es deliberada y se considera apropiada para el estado actual de madurez de
este tipo de sistemas.

Debe reconocerse asimismo una limitación inherente al enfoque de evaluación adoptado: los
conjuntos de datos empleados provienen de laboratorios y de campañas concretas, por lo que
los resultados obtenidos no son directamente extrapolables a la diversidad de un entorno de
producción real. Esta limitación se explicita en el análisis de resultados y en las
conclusiones.

\section{Licencias y propiedad intelectual}

El trabajo se apoya íntegramente en componentes de código abierto y conjuntos de datos
públicos, cuyas condiciones de uso se han verificado y respetado. Los conjuntos de datos
empleados se distribuyen bajo licencias que permiten su uso con fines de investigación, y las
bases de conocimiento y formatos utilizados —MITRE ATT\&CK, OCSF, STIX, Sigma— son recursos
abiertos publicados para su adopción por la comunidad.

El Anexo~\ref{anx:licencias} detalla la licencia de cada componente, conjunto de datos y
base de conocimiento. Dos casos merecen mención: el bus de eventos se distribuye bajo una
licencia de código fuente disponible que permite su uso gratuito salvo para ofrecerlo como
servicio comercial, y el modelo de lenguaje previsto para la extracción de técnicas se
publica bajo una licencia abierta con restricciones de uso responsable, compatibles con la
finalidad defensiva del trabajo.

\section{Impacto social y económico}

El impacto potencial más directo del trabajo se sitúa en el ámbito de las organizaciones que
carecen de capacidad para operar un centro de operaciones de seguridad dotado. Las
herramientas comerciales que proporcionan capacidades equivalentes tienen un coste que las
sitúa fuera del alcance de una parte considerable del tejido empresarial. Una arquitectura
construida íntegramente sobre componentes abiertos, desplegable sobre recursos modestos y
que reduce el volumen de alertas a revisar, reduce la barrera de entrada a este tipo de
capacidades.

Debe señalarse, no obstante, que la reducción de la barrera tecnológica no elimina la
necesidad de criterio experto: el sistema ordena y contextualiza, pero la interpretación y la
respuesta siguen requiriendo conocimiento especializado.

\section{Impacto ambiental}

El procesamiento continuo de telemetría tiene un coste energético asociado. Varias decisiones
de diseño adoptadas por motivos de eficiencia operativa tienen también efecto en este plano:
la selección de componentes de bajo consumo de recursos, la organización del despliegue en
perfiles que evitan mantener servicios activos innecesariamente, la aplicación de
codificaciones de compresión específicas por columna en el almacén analítico, y la
restricción del aprendizaje automático a modelos pequeños entrenados en la propia estación
de trabajo, descartando el entrenamiento de modelos de gran tamaño.

Esta última decisión, adoptada inicialmente por limitación de recursos, resulta relevante
también desde esta perspectiva: el coste energético del entrenamiento de modelos grandes es
considerable, y el trabajo demuestra que el problema abordado no lo requiere.

El desarrollo se ha realizado íntegramente sobre una única estación de trabajo, sin recurso a
infraestructura de centro de datos.

\chapter{Presupuesto económico}
\label{cap:presupuesto}

Este anexo estima el coste de realizar el trabajo. Se estructura en costes de personal,
de material y de licencias. Las cifras de personal y de amortización se cerrarán al final del
desarrollo, cuando el cómputo de horas dedicadas sea completo.

\section{Costes de personal}

El coste de personal se calculará como el producto de las horas dedicadas a cada fase por un
coste horario de referencia para un ingeniero recién titulado. Las horas se obtienen del
registro de trabajo del proyecto, que anota cada sesión con su objetivo y su resultado.

% TODO: tabla de horas por fase y coste horario de referencia.

\section{Costes de material}

Todo el desarrollo se ha realizado sobre una única estación de trabajo propia, con procesador
de 16 núcleos, 64~GB de memoria y GPU dedicada, sin infraestructura de nube ni de centro de
datos. Su coste se imputará como amortización lineal durante los meses de uso, sobre una vida
útil de referencia.

% TODO: precio de adquisición, vida útil y meses de uso.

\section{Costes de licencias}

El coste de licencias es \textbf{nulo}. Todos los componentes empleados son de código abierto
o de uso gratuito (Anexo~\ref{anx:licencias}), y los conjuntos de datos y las bases de
conocimiento son públicos. Las plataformas comerciales que ofrecen capacidades equivalentes,
como los SIEM con correlación y priorización de incidentes, se licencian habitualmente por
volumen de datos ingeridos o por equipo supervisado, lo que constituye precisamente la barrera
de entrada que el trabajo pretende reducir.

% \section{Resumen del presupuesto}
% TODO: tabla resumen (personal + material + licencias, costes generales, beneficio e IVA).

\chapter{Esquema de campaña}
\label{anx:esquema}

Las campañas de todos los patrones del Pilar~A se emiten como hallazgos \textbf{OCSF 1.5.0
«Detection Finding»} (\texttt{class\_uid} 2004, \texttt{type\_uid} 200401) con una extensión
propia, \texttt{tfg}. Se publican en el \textit{topic} \texttt{det.campaigns}, sin caducidad,
y se guardan en la tabla \texttt{tfg.campaigns} del almacén analítico.

\section{Campos del estándar}

\begin{table}[H]
\centering
\small
\begin{tabular}{|l|p{8.2cm}|}
\hline
\textbf{Campo} & \textbf{Contenido} \\
\hline
\texttt{start\_time}, \texttt{end\_time}, \texttt{duration} & Inicio, fin y duración de la campaña, en milisegundos \\
\texttt{time} & Igual a \texttt{end\_time}, para que reprocesar produzca el mismo documento \\
\texttt{severity\_id} & Máxima gravedad de las reglas que aportan evidencia \\
\texttt{confidence\_score} & Suma de los factores de confianza, de 0 a 100 \\
\texttt{resources[]} & Equipos implicados \\
\texttt{finding\_info.uid} & Identidad de la campaña \\
\texttt{finding\_info.analytic} & Patrón que la reconoce (P1, P2, P3) \\
\texttt{finding\_info.attacks[]} & Técnicas observadas en la evidencia, sin táctica asociada \\
\texttt{finding\_info.related\_events[]} & Alertas y eventos de evidencia, una vez cada uno \\
\hline
\end{tabular}
\caption{Campos del estándar OCSF empleados.}
\end{table}

\section{Extensión propia}

\begin{table}[H]
\centering
\small
\begin{tabular}{|l|p{9cm}|}
\hline
\textbf{Campo} & \textbf{Contenido} \\
\hline
\texttt{tfg.pattern\_id}, \texttt{tfg.dataset\_id} & Patrón y conjunto de datos de origen \\
\texttt{tfg.devices[]} & Equipos con su papel: \textit{host} en P1 y P3, origen y destino en P2 \\
\texttt{tfg.tactic\_names[]} & Tácticas de la evidencia, por nombre \\
\texttt{tfg.stages[]} & Etapas ordenadas, cada una con su intervalo, sus técnicas y sus evidencias \\
\texttt{tfg.confidence\_factors} & Factores presentes y su peso; su suma es la confianza \\
\texttt{tfg.metrics} & Métricas propias del patrón (aristas, transiciones, alertas en la ventana) \\
\hline
\end{tabular}
\caption{Campos de la extensión \texttt{tfg}.}
\end{table}

\section{Reglas de construcción}

\begin{itemize}
    \item \textbf{Identidad}: \texttt{md5(dataset | patrón | clave | inicio en ms)}, donde la
    clave es el equipo en P1 y P3 y el par origen--destino en P2. Es determinista.
    \item \textbf{Orden de las etapas}: por el segundo de la primera evidencia de cada etapa y,
    dentro del mismo segundo, por el orden lógico del patrón.
    \item \textbf{Una evidencia puede probar dos etapas} (el acceso remoto a LSASS mediante
    WinRM es a la vez ejecución remota y acceso a credenciales). Aparece en las dos etapas,
    pero una sola vez en \texttt{related\_events}.
    \item \textbf{Validación}: el proceso de volcado comprueba la clase, la coherencia de los
    tiempos, que la confianza sea la suma de sus factores, que las etapas estén numeradas de 1
    a $n$ con su evidencia y que no haya evidencias repetidas.
\end{itemize}

\section{Ejemplo}

El Listado~\ref{lst:campana-p2} muestra, abreviada, la campaña de movimiento lateral del
primer día de APT29.

\begin{lstlisting}[caption={Campaña P2 de APT29, abreviada.},label={lst:campana-p2}]
{
  "class_uid": 2004, "type_uid": 200401, "activity_id": 1,
  "start_time": 1588388963801, "end_time": 1588389479493,
  "duration": 515692, "severity_id": 4, "confidence_score": 100,
  "message": "Movimiento lateral SCRANTON -> NASHUA",
  "finding_info": {
    "uid": "07c1426c17179d180b24073e097edfdc",
    "analytic": {"uid": "P2", "name": "P2_lateral_movement"},
    "attacks": [{"technique": {"uid": "T1021"},
                 "sub_technique": {"uid": "T1021.006"}}],
    "related_events_count": 8
  },
  "tfg": {
    "pattern_id": "P2",
    "devices": [{"hostname": "SCRANTON", "role": "origen"},
                {"hostname": "NASHUA", "role": "destino"}],
    "stages": [{"order": 1, "name": "conexion_lateral", "count": 17}],
    "confidence_factors": {"arista": 0.30, "bilateral": 0.15,
      "proceso_interactivo": 0.10, "ejecucion_remota": 0.20,
      "servicio": 0.15, "credenciales": 0.10},
    "metrics": {"aristas": 17, "transiciones": 13,
      "alertas_absorbidas": 115, "retraso_max_ms": 1004}
  }
}
\end{lstlisting}

\chapter{Implementación de los patrones en Flink SQL}
\label{anx:patrones}

Cada patrón se organiza en los mismos pasos, cada uno en un fichero SQL que se ejecuta con el
cliente de Flink y una definición común de tablas y vistas que se carga al inicio:

\begin{enumerate}
    \item \textbf{Flujo}: publica en un \textit{topic} intermedio
    (\texttt{det.p1\_flujo}, \texttt{det.p2\_flujo}, \texttt{det.p3\_flujo}) la evidencia del
    patrón, sin repeticiones y en orden de tiempo de evento.
    \item \textbf{Campañas}: correlación en dos niveles sobre ese flujo.
    \item \textbf{Emisión}: construcción del hallazgo OCSF y publicación en
    \texttt{det.campaigns}.
    \item \textbf{Flujo continuo}: el mismo patrón en modo \textit{streaming}, con una ventana
    de sesión como segundo nivel, frente a su resultado por lotes.
\end{enumerate}

El Listado~\ref{lst:p1} reproduce los dos niveles de P1 en su versión por lotes.

\begin{lstlisting}[language=SQL,caption={P1: transiciones (nivel 1) y campañas (nivel 2).},label={lst:p1}]
-- NIVEL 1: dos evidencias consecutivas del mismo equipo a <= 10 min,
-- al menos una de credenciales (transicion bidireccional).
CREATE TEMPORARY VIEW transiciones_p1 AS
SELECT *
FROM (SELECT dataset, host, ts, `time` AS t_ms, cred FROM p1_flujo)
MATCH_RECOGNIZE (
    PARTITION BY dataset, host
    ORDER BY ts
    MEASURES E.t_ms AS t_ini_ms, C.t_ms AS t_fin_ms
    ONE ROW PER MATCH
    AFTER MATCH SKIP TO NEXT ROW
    PATTERN (E C)
    DEFINE
        C AS (C.cred = 1 OR E.cred = 1)
         AND TIMESTAMPDIFF(SECOND, E.ts, C.ts) <= 600
) AS t;

-- NIVEL 2: un hueco de mas de 10 min entre transiciones abre campana.
CREATE TEMPORARY VIEW transiciones_marcadas_p1 AS
SELECT *,
       CASE WHEN LAG(t_fin_ms) OVER (PARTITION BY dataset, host
                                      ORDER BY t_ini_ms) IS NULL
              OR (t_ini_ms - LAG(t_fin_ms) OVER (PARTITION BY dataset,
                                      host ORDER BY t_ini_ms)) > 600000
            THEN 1 ELSE 0 END AS abre
FROM transiciones_p1;

CREATE TEMPORARY VIEW campanas_base_p1 AS
SELECT dataset, host, campana,
       MIN(t_ini_ms) AS inicio, MAX(t_fin_ms) AS fin,
       COUNT(*) AS transiciones
FROM (SELECT *, SUM(abre) OVER (PARTITION BY dataset, host
             ORDER BY t_ini_ms ROWS BETWEEN UNBOUNDED PRECEDING
             AND CURRENT ROW) AS campana
      FROM transiciones_marcadas_p1)
GROUP BY dataset, host, campana;
\end{lstlisting}

En modo continuo, el segundo nivel se sustituye por una ventana de sesión sobre el tiempo de
cada transición, que Flink proporciona como \texttt{MATCH\_ROWTIME()}:

\begin{lstlisting}[language=SQL,caption={P3: segundo nivel en modo continuo.},label={lst:p3-sesion}]
SELECT dataset, host,
       MIN(t_ini_ms) AS inicio, MAX(t_fin_ms) AS fin,
       COUNT(*) AS transiciones
FROM TABLE(SESSION(TABLE transiciones_p3_rt PARTITION BY clave,
                   DESCRIPTOR(rt), INTERVAL '10' MINUTE))
GROUP BY clave, dataset, host, window_start, window_end
HAVING SUM(con_otra) >= 1;
\end{lstlisting}

\section{Restricciones de Flink 2.2 encontradas}

\begin{itemize}
    \item No se admiten cuantificadores voraces al final de un patrón, lo que obliga a la
    correlación en dos niveles (Sección~\ref{sec:dos-niveles}).
    \item \texttt{DEFINE} no admite subconsultas y \texttt{MEASURES} no admite
    \texttt{COUNT(DISTINCT)}.
    \item \texttt{WITHIN} limita la duración total de la coincidencia, no el hueco entre
    elementos.
    \item \texttt{LAG} no acepta marcas de tiempo con zona; se opera sobre milisegundos.
    \item Una ventana de sesión sobre una vista de una tabla cuyo tiempo es una columna
    calculada pierde el atributo de tiempo; se aplica sobre la tabla.
    \item Una vista que contiene una ventana de sesión con partición de varias columnas no se
    puede volver a expandir al consultarla; se escribe como subconsulta.
    \item \texttt{JSON\_OBJECT} inserta como texto los objetos que llegan de una subconsulta;
    el documento de campaña se compone con \texttt{LISTAGG} y \texttt{CONCAT}.
    \item Los literales de texto en un \texttt{CASE} se rellenan con espacios hasta la mayor
    longitud; se convierten explícitamente a \texttt{STRING}.
\end{itemize}

\chapter{Entorno y reproducibilidad}
\label{anx:reproducibilidad}

\section{Servicios}

Todos los servicios se despliegan con Docker Compose sobre un único anfitrión y se agrupan en
perfiles, de modo que solo se levanta lo necesario para cada tarea.

\begin{table}[H]
\centering
\small
\begin{tabular}{|l|l|l|l|}
\hline
\textbf{Servicio} & \textbf{Versión} & \textbf{Perfil} & \textbf{Estado} \\
\hline
Redpanda y su consola & 25.2.1 / 3.1.0 & \texttt{core} & En uso \\
Apache Flink (gestor y trabajador) & 2.2.1 & \texttt{stream} & En uso \\
ClickHouse & 25.3 & \texttt{store} & En uso \\
Prometheus y Grafana & 3.1.0 / 11.5.1 & \texttt{observability} & En uso \\
MLflow & 2.19.0 & \texttt{mlops} & En uso \\
Portal web & --- & \texttt{portal} & En uso (demostración) \\
Neo4j con GDS & --- & \texttt{graph} & Fase 5 \\
OpenCTI & --- & \texttt{cti} & Fase 4 \\
Dagster & --- & \texttt{orchestration} & Fase 4 \\
FastAPI & --- & \texttt{api} & Fase 6 \\
\hline
\end{tabular}
\caption{Servicios del sistema y perfiles de despliegue.}
\label{tab:servicios}
\end{table}

\section{Topics del bus de eventos}

\begin{table}[H]
\centering
\small
\begin{tabular}{|l|c|l|}
\hline
\textbf{Topic} & \textbf{Particiones} & \textbf{Contenido} \\
\hline
\texttt{raw.host.windows} & 3 & Telemetría cruda de los equipos \\
\texttt{norm.events} & 6 & Eventos normalizados a OCSF \\
\texttt{det.alerts} & 3 & Alertas de las reglas Sigma \\
\texttt{det.p1\_flujo}, \texttt{det.p2\_flujo}, \texttt{det.p3\_flujo} & 3 & Evidencia de cada patrón, en orden de tiempo \\
\texttt{det.campaigns} & 3 & Campañas en el esquema común \\
\texttt{dlq.*} & 1 & Rechazos de cada etapa, con su motivo \\
\hline
\end{tabular}
\caption{Topics empleados en el Pilar A.}
\end{table}

La partición de \texttt{norm.events} se decide por la identidad del evento (conjunto de
datos, equipo, canal y número de registro), que reparte la carga de forma uniforme entre los
detectores en paralelo. Los \textit{topics} del corpus se crean sin caducidad: la retención
por defecto de siete días los vació en una ocasión sin producir ningún error.

\section{Reproducción del corpus}

El corpus de evaluación se carga con un único guion, que recrea los \textit{topics},
reproduce los nueve conjuntos de datos y los procesa con grupos de consumo nuevos:

\begin{lstlisting}[language=bash,caption={Carga del corpus del Pilar A.},label={lst:carga}]
docker compose --profile core --profile stream --profile store up -d
DETECT_INSTANCIAS=6 ./scripts/load_f3_corpus.sh
\end{lstlisting}

Las cifras esperadas tras la carga son las de la Tabla~\ref{tab:corpus-f3}. Como los
identificadores de eventos, alertas y campañas son deterministas y las tablas del almacén
deduplican por ellos, la carga puede repetirse o reanudarse sin duplicar resultados.

\chapter{Licencias de componentes y datos}
\label{anx:licencias}

\begin{table}[H]
\centering
\footnotesize
\begin{tabularx}{\textwidth}{|l|l|X|}
\hline
\textbf{Componente o recurso} & \textbf{Licencia} & \textbf{Uso en el trabajo} \\
\hline
\multicolumn{3}{|l|}{\textit{Datos y bases de conocimiento}} \\
\hline
OTRF Security-Datasets (incluida APT29) & MIT & Corpus de evaluación \\
MITRE ATT\&CK v19.2 & Términos de uso de MITRE & Catálogo de técnicas y referencia común \\
ThreatFox (abuse.ch) & Uso gratuito (\textit{fair use}) & Feed de indicadores \\
AnnoCTR & CC BY-SA 4.0 & Entrenamiento y evaluación de ML-2 \\
Informes de Microsoft Security Blog & Derechos de Microsoft (copia local) & Evaluación de ML-2 sobre APT29 \\
Reglas Sigma (SigmaHQ) & Detection Rule License 1.1 & Reglas de detección \\
Esquema OCSF & Apache 2.0 & Esquema canónico \\
\hline
\multicolumn{3}{|l|}{\textit{Componentes en uso}} \\
\hline
Docker Engine & Apache 2.0 & Contenedores \\
Redpanda & Business Source License 1.1 & Bus de eventos \\
Apache Flink & Apache 2.0 & Procesamiento y CEP \\
ClickHouse & Apache 2.0 & Almacén analítico \\
pySigma & LGPL 2.1 & Traducción de reglas \\
scikit-learn & BSD de 3 cláusulas & Detección de anomalías \\
MLflow & Apache 2.0 & Registro de experimentos \\
Prometheus & Apache 2.0 & Métricas \\
Grafana & AGPL 3.0 & Paneles \\
Sysmon & Sysinternals (gratuita) & Telemetría de Windows \\
Atomic Red Team & MIT & Laboratorio de emulación \\
VirtualBox & GPL 3.0 & Laboratorio de emulación \\
React & MIT & Portal web \\
Neo4j Community, GDS y APOC & GPL 3.0 (APOC: Apache 2.0) & Grafo de conocimiento y de ataque \\
OpenCTI & Apache 2.0 & Comparativa de plataformas CTI \\
PyTorch & BSD de 3 cláusulas & Afinado de modelos \\
Transformers (Hugging Face) & Apache 2.0 & Afinado de modelos \\
SecureBERT 2.0 y ModernBERT-base & Apache 2.0 & Extracción de técnicas (ML-2) \\
\hline
\multicolumn{3}{|l|}{\textit{Componentes previstos}} \\
\hline
Dagster & Apache 2.0 & Orquestación \\
Soufflé & UPL 1.0 & Razonamiento Datalog \\
FastAPI & MIT & API \\
SHAP & MIT & Explicación de la puntuación \\
MinIO & AGPL 3.0 & Almacén de modelos \\
Kubernetes & Apache 2.0 & Despliegue \\
CALDERA & Apache 2.0 & Simulación de campañas \\
\hline
\end{tabularx}
\caption{Licencias de los componentes, datos y bases de conocimiento.}
\label{tab:licencias}
\end{table}

Ninguna de estas licencias impone coste ni impide el uso académico. Los informes de Microsoft
no se redistribuyen: solo se usan copias locales y se citan fragmentos breves. Las licencias de tipo
\textit{copyleft} (GPL y AGPL) solo obligarían a publicar el código fuente si se
distribuyera una versión modificada del componente, lo que no ocurre en este trabajo.


\end{document}
